% \iffalse
%
% Copyright 2010 Stephen Hicks, All rights reserved.
% marginfix.dtx - 29 Jul 2010
% originally floatprobsbegone, created 21 Mar 2008
%
% Run LaTeX on this document to produce documentation.
% Run LaTeX on marginfix.ins to produce the package.
%<*driver>
\ProvidesFile{marginfix.dtx}
%</driver>
%
%<package>\NeedsTeXFormat{LaTeX2e}
%<package>\ProvidesPackage{marginfix}%
           [2026/10/02 v1.5 Fix Margin Paragraphs]
%<*driver>
\documentclass{ltxdoc}
\CheckSum{3193}
%\OnlyDescription % (un)comment this line to show (hide) source code
\RecordChanges
\EnableCrossrefs
\CodelineIndex   % (un)comment this line to index source by page (line)
\begin{document}
  \newcommand*\Lopt[1]{\textsf {#1}}
  \newcommand*\lit[1]{\texttt{\char`#1}} %% literal character (funny catcodes)
  \parindent0pt
  \def\*#1{\texttt{\string#1}} %% sdh - |...| doesn't work in headings
  \makeatletter
  %
  \let\pkg\textsf
  \newcount\mac@depth\mac@depth\z@
  \newcommand\@macros{}\newcommand\@endmacros{}
  \catcode`&3  %% we use a funny catcode to ensure never used.
  \def\@macros#1,{\macro{#1}\global\advance\mac@depth\@ne\relax
    \@ifnextchar&\@gobble\@macros}
  \def\@endmacros{\let\mac@next\relax\ifnum\mac@depth>\z@
    \endmacro\let\mac@next\@endmacros
    \global\advance\mac@depth\m@ne\fi\mac@next}
  \newenvironment{macros}[1]{\@macros#1,&}{\@endmacros}
  \catcode`&4  %% put it back 
  \makeatother %% must be balanced for character table to work properly
  %
  \DocInput{marginfix.dtx}
  \setcounter{IndexColumns}{2}
  \PrintIndex
  \PrintChanges
\end{document}
%</driver>
% \fi
% \changes{v0.0}{2008/03/21}
%       {(SDH) Initial version of floatprobsbegone.}
% \changes{v0.9}{2010/08/18}
%       {(SDH) Initial CTAN version.}
% \changes{v0.9.1}{2010/08/28}
%       {(SDH) Fix bug where we alternated sides in article.}
% \changes{v1.0}{2013/09/01}
%       {(Dario Buttari) Alternate order of \cs{marginpar} arguments to be
%        consistent with the original macro.}
% \changes{v1.0}{2013/09/01}
%       {(Dario Buttari) Fix bug in deferred note spacing.}
% \changes{v1.0}{2013/09/01}
%       {(Dario Buttari) Fix issues with \cs{marginheightadjustment} and
%        \cs{extendmargin} not being applied and reset consistently.}
% \changes{v1.0}{2013/09/01}
%       {(SDH) Stop gobbling vertical stretch on page.  Margin notes will
%        not line up properly if page is stretched.}
% \changes{v1.0}{2013/09/01}
%       {(SDH) Add margin phantoms.}
% \changes{v1.1}{2013/09/08}
%       {(SDH) Globally calculate margin phantoms over 4 passes.}
% \changes{v1.1}{2013/09/08}
%       {(SDH) Add \cs{topskip} to notes at the top of the margin.}
% \changes{v1.2}{2020/05/06}
%       {(SDH) Fix long-standing bug where margin notes called out in
%        the last few points of a page were being entirely dropped.}
% \changes{v1.3}{2026/06/23}
%       {(Daniel Bosk) Keep each margin note on the page of its callout:
%        record callout pages in the .aux and, on the next run, hold back
%        notes that \LaTeX{} would otherwise set on the previous page.}
% \changes{v1.4}{2026/09/22}
%       {(Daniel Bosk) Hook into the kernel's \cs{@outputbox@attachfloats},
%        so that notes are attached under the standard classes with
%        \LaTeX{} 2020-10-01 or later (the kernel's first aid only did this
%        for v1.2).}
% \changes{v1.5}{2026/09/27}
%       {(Daniel Bosk) Add deferrable notes (\cs{marginnotedeferrable}) and
%        \cs{marginnotepriority}: warn when deferrable notes crowd an
%        ordinary note off its page, and optionally move or split them
%        to keep the ordinary notes with their callouts; split notes
%        between paragraphs where possible (\cs{marginsplitslack}), with
%        the cues \cs{marginnotecontinued} and \cs{marginnotecontinues}.}
% \changes{v1.5}{2026/10/02}
%       {(Daniel Bosk) Add \cs{marginnotebesidefloats}: optionally set
%        notes beside bottom floats that stay in the text block, unless
%        a float is wider than the column, has a side caption, or says
%        \cs{marginnotereservefloat}.}
%
% \GetFileInfo{marginfix.dtx}
% \title{\Lopt{marginfix} package documentation}
% \author{Stephen Hicks\\%
%    \texttt{sdh33@cornell.edu}\\%
%    \texttt{https://github.com/shicks/marginfix}}
% \date{\fileversion{} -- \filedate}
% \maketitle
%
% \part*{Usage}
% \section{Overview}
% Authors using \LaTeX\ to typeset books with significant margin material
% often run into the problem of long notes running off the bottom of the
% page.  A typical workaround is to insert \cs{vshift}s by hand, but
% this is a tedious process that is invalidated when pagination changes.
% Another workaround is \pkg{memoir}'s \cs{sidebar} function, but this can be
% unsatisfying for short textual notes, and standard marginpars cannot
% be mixed with sidebars.  This package implements a solution to
% make marginpars "just work" by keeping a list of floating inserts and
% arranging them intelligently in the output routine.  The credit for the
% concept behind this algorithm goes to Prof.\ Andy Ruina, who employed me
% to work on some of his textbook macros in 2007--9.
%
% As of v1.3 the package also keeps each note on the page of its callout.
% A \cs{marginpar} is a float, so \LaTeX{} can attach it to the page it is
% assembling even when the page then breaks \emph{before} the line that
% called it---leaving a note whose mark is at the top of a page set in the
% margin of the previous page.  marginfix now records each callout's page in
% the \texttt{.aux} file and, on the next run, holds such a note back until
% its own page is shipped (requesting a rerun through the usual ``Label(s)
% may have changed'' mechanism).  This is automatic and needs no markup;
% like cross-references it simply takes an extra pass to settle.
%
% As of v1.5 notes can be marked as \emph{deferrable}: commentary that may
% move to a later page, or be split across pages, so that the ordinary notes
% after it (for instance footnotes set in the margin) stay beside their
% callouts.  See section~\ref{sec:priority}.
%
% \section{Options}
% There are currently no options that do anything yet.
%
% \section{Commands}
% For the most part, this is a drop-in replacement.  Simply include
% a call to |\usepackage{marginfix}| to the preamble, use \cs{marginpar}
% normally and hope for the best.
% In the event, however, that it doesn't work exactly as hoped,
% there are a number of tweaks that the user can apply.
%
% \DescribeMacro{\marginskip}
% Calling \cs{marginskip}\marg{length} will insert an incompressible
% skip in the margin.  These skips will force neighboring notes on
% the same page to be separated, but will disappear at the top or
% bottom of a margin.
%
% \DescribeMacro{\clearmargin}\DescribeMacro{\softclearmargin}
% In an analog to \cs{clearpage}, \cs{clearmargin} prevents any further
% material from being added to the current margin.  These calls are
% cumulative, so that two \cs{clearmargin}s in a row will produce a
% completely empty margin on the next page as well.  If this is not
% the desired effect, use \cs{softclearmargin}, which is effectively
% idempotent: multiple calls have the same effect as one call to
% end the current margin.
%
% \DescribeMacro{\extendmargin}
% If a page has too much margin material to fit and an important note
% is floating to the next page, \cs{extendmargin}\marg{length} will
% extend the margin (for the current page only) by the given length.
% If the length is negative, the margin will shrink.  Multiple calls
% on the same page are cumulative.
%
% \DescribeMacro{\mparshift}
% To adjust the position of a single note, use \cs{mparshift}\marg{length}
% before a call to \cs{marginpar}.  Positive lengths move it down the page.
% This essentially shifts the call-out location, so the actual position of
% the note might not change if the margin is sufficiently crowded.
% Multiple calls before the same note are cumulative.
%
% \DescribeMacro{\marginheightadjustment}
% If all the margins are the wrong size, the height of the margin on every
% page can be adjusted by assigning a non-zero value to the dimension
% register \cs{marginheightadjustment} (as in 
% \cs{marginheightadjustment}\texttt{=}$\langle$\emph{length}$\rangle$).
% This is effectively the same as a call to
% \cs{extendmargin} on every page.
%
% \DescribeMacro{\marginposadjustment}
% Similarly, if all the margin notes are in the wrong place, the callout
% positions can be adjusted globally by assigning a non-zero value to the
% dimension register \cs{marginposadjustment}.  This is effectively the same
% as a call to \cs{mparshift} before every note.  This is particularly useful
% at present because the height of the line on which the margin note is called
% is currently only estimated, and appears to be off by a point or two.
% This may get fixed in the future, but until then, the adjustment is
% possibly the easiest workaround.
%
% \DescribeMacro{\blockmargin}\DescribeMacro{\unblockmargin}
% As of version 1.0, we now support ``margin phantoms'': sections of the
% margin in which no notes will be placed, which can be useful for large
% figures that jut into the margin (note: margins already move out of the
% way of floats, regardless of whether or not they extend into the margin;
% this is mainly for in-place figures or equations).  The easiest way to
% block off part of the margin is to call \cs{blockmargin} before the
% extended content and \cs{unblockmargin} afterwards.  No margin notes
% will be placed between these two points (though one must be careful:
% if one of these is called in horizontal mode, the toggle will occur at
% the \emph{top} of the current line).  Each of these commands takes an
% optional argument: \cs{blockmargin}\oarg{pos} will begin the margin
% block at a position \emph{pos} below the current position (or above
% if \emph{pos} is negative), and \cs{unblockmargin}\oarg{pos} will
% likewise end the block at a position \emph{pos} below the current
% position.
%
% \DescribeMacro{\marginphantom}
% Margin phantoms may also be called out in place with a known size
% using \cs{marginphantom}\oarg{pos}\marg{size}, which is essentially
% equivalent to \cs{blockmargin}\oarg{pos}\cs{unblockmargin}\oarg{pos$+$size}.
% Either argument may be negative to refer upward rather than downward.
%
% \section{Deferrable notes and priority}
% \label{sec:priority}
% marginfix sets the notes of a page in the order of their callouts, and when
% the margin is full, the first note that doesn't fit moves to the next page,
% taking every later note with it.  That is usually what we want, but not
% when the notes are of two kinds.  Take a book with its footnotes in the
% margin (\pkg{memoir}'s \cs{footnotesinmargin}) and long commentary notes
% in the same margin, say for the teacher's edition.  A few commentary notes
% called out high on a page can fill the margin, and then a footnote called
% out lower down moves to the next page, away from its mark, while the
% commentary, which could just as well have continued on the next page,
% stays.
%
% \DescribeMacro{\marginnotedeferrable}
% Writing \cs{marginnotedeferrable} in the text of a note marks the note as
% \emph{deferrable}:
% \begin{verbatim}
% \marginpar{\footnotesize\marginnotedeferrable Why this example: ...}
% \end{verbatim}
% All other notes are \emph{ordinary}.  Put \cs{marginnotedeferrable} after
% the note's font commands, since it records the font for the cues of a
% split note (see below).  A package that makes commentary notes can mark
% them all by putting \cs{marginnotedeferrable} in the note text it builds,
% guarded by |\ifdefined\marginnotedeferrable| so that it still works with
% older versions of marginfix.  Used outside a margin note,
% \cs{marginnotedeferrable} does nothing but warn.
%
% \DescribeMacro{\marginnotepriority}
% What marginfix does with the deferrable notes is set by
% \cs{marginnotepriority}\marg{mode}, in the preamble or anywhere in the
% document.  The modes are:
% \begin{description}
% \item[\texttt{warn}] (the default) Notes are set exactly as without
%   deferrable notes, but whenever an ordinary note moves to a later page
%   while deferrable notes ahead of it stayed, or didn't fit either, we warn,
%   naming the ordinary note and the deferrable ones by the file and line
%   where each note ends.  If a deferrable note that didn't fit took the
%   notes after it along, the warning says so.  Then you know which note to
%   shorten or re-anchor.
% \item[\texttt{defer}] A deferrable note that is in the way moves, whole, to
%   the next page, so that the ordinary notes after it fit on their page.
%   The notes of each kind keep their order.  On the next page, the moved
%   note comes after that page's ordinary notes, up to its first deferrable
%   note, so that it doesn't push the next page's footnotes away from their
%   marks in turn.  We warn about every note that moves this way.
%   A deferrable note that doesn't fit at all moves to the next page too,
%   but no longer takes the notes after it along: the ordinary notes after
%   it that fit are still set on the page.  The deferrable notes after it
%   move with it, though, so that they keep their order; and an ordinary
%   note that doesn't fit still takes every later note along, so that
%   footnotes never trade places.
% \item[\texttt{split}] Like \texttt{defer}, but the note that is in the way
%   is split, and so is a deferrable note that doesn't fit: what fits stays
%   on the page, where its callout is, and the rest continues at the top of
%   the next page.  If the note can't be split (too
%   little room left, nothing to break it at, or a margin interrupted by
%   \cs{blockmargin}), it moves whole as with \texttt{defer}.  We warn about
%   every note that is split or moved.  See below for where notes are split
%   and how the parts are marked.
% \item[\texttt{off}] As \texttt{warn}, but without the warnings.
% \end{description}
% An ordinary note can still move to a later page, when the ordinary notes
% alone don't fit; nothing can help that.  A deferrable note may move more
% than one page if the following margins are full as well.
%
% The mode is used when a page is built, and \TeX{} builds a page some
% time after it has read the text on it.  So a new mode set in the middle of
% the document applies from the next page built after the command is read,
% which may be the page the command is on; put a \cs{clearpage} before it
% to be sure.
%
% \DescribeMacro{\marginsplitslack}
% In the mode \texttt{split}, a note is split between paragraphs if it can
% be: at the last paragraph break where the first part fits, unless that
% leaves more than \cs{marginsplitslack} of the room unused (or more than
% half of it).  The default is four lines of the text.  Otherwise the note
% is split between lines, where the
% penalties of the note's paragraphs are lowest, so that, with the usual
% club and widow penalties, no single line of a paragraph is left on either
% side.  Only when no break fits that bill do we take the last line that
% fits.  A note is never split in the middle of a line, but it may be split
% in the middle of a sentence.
%
% \DescribeMacro{\marginnotecontinued}\DescribeMacro{\marginnotecontinues}
% The rest of a split note starts with a line saying
% \cs{marginnotecontinued}\marg{text}, by default ``(continued)'', so that
% it doesn't read as a note of its own.  Its first part can end with a line
% saying \cs{marginnotecontinues}\marg{text}, set flush right, for
% instance an arrow; by default that is empty, and then there is no such
% line.  (A line at the end of the first part takes room on the crowded
% page, while the rest starts at the top of the next margin, where the
% reader finds it anyway.)  An empty \meta{text} turns a cue off.  The cues
% are set in the font of the note where it said \cs{marginnotedeferrable},
% so put \cs{marginnotedeferrable} after the note's font commands.  A
% package can set the cues for a language, e.g.\
% |\marginnotecontinued{(forts.)}| for Swedish.
%
% Like the mode, these settings are used when a page is built, not where
% the note is, so they are all global: \cs{marginnotepriority},
% \cs{marginnotecontinued} and \cs{marginnotecontinues} take effect
% whatever group they are in, from the next page built.
% \cs{marginsplitslack} is a length, and the value that counts is the one
% in force when the page is built.  Set it in the preamble with
% \cs{setlength}, or anywhere with |\global\marginsplitslack=|\meta{length};
% a \cs{setlength} in a group is forgotten at its end and may or may not be
% in force when the page is built.
%
% \section{Notes beside floats}
% \label{sec:floats}
% marginfix sets notes only in the margin beside the text of a page, not
% beside the floats at its top or bottom, since a float may extend into the
% margin (a side caption, or a figure wider than the text).  For a plain
% table at the bottom of a page that leaves a stretch of margin empty
% while notes called out on the page move on to the next.
%
% \DescribeMacro{\marginnotebesidefloats}
% \cs{marginnotebesidefloats}\marg{mode} decides what happens beside the
% floats at the bottom of a page.  With \texttt{reserve}, the default, the
% margin beside them stays empty, as it always has.  With \texttt{use}, the
% notes may continue beside them, unless one of them extends into the
% margin; then the margin beside all the bottom floats of that page stays
% empty.  A float counts as extending into the margin if
% \begin{itemize}
% \item any of its lines is wider than the column, or overfull (by more
%   than 1\,pt): a table wider than the text, or text set wider on either
%   side, as with \pkg{memoir}'s \texttt{adjustwidth} with a negative
%   margin, which also catches text that sticks out on the left;
% \item it uses one of \pkg{memoir}'s side captions or legends
%   (\texttt{sidecaption}, \texttt{sidecontcaption}, \texttt{sidelegend},
%   \texttt{sidenamedlegend}); or
% \item it says \cs{marginnotereservefloat}, see below.
% \end{itemize}
% Floats at the top of a page always keep the margin beside them empty.
% Like the other settings, the mode is global, and used from the next
% page built.
%
% \DescribeMacro{\marginnotereservefloat}
% marginfix looks at the lines of a float, and at boxes at the end of a
% line (a minipage, \cs{parbox} or table placed with \texttt{t} or
% \texttt{b}), but not further.  It doesn't look into a minipage, \cs{parbox}
% or table centred on its line (the default), the inside of an \cs{fbox}, a
% formula, or a box followed by text on its line.  And \TeX{} can't see ink
% outside the boxes: a \cs{rlap} or \cs{llap}, a \cs{makebox} that lets
% its content stick out, negative space that pulls text out of its line, a
% TikZ overlay, or a picture drawn beyond its bounding box.  A float in
% which any of these sticks out into the margin goes unnoticed; write
% \cs{marginnotereservefloat} in it to keep the margin beside it empty.
% It applies to the float it is in, and only does anything with
% \cs{marginnotebesidefloats}\texttt{\{use\}}; outside a float it warns.
% (Used where \cs{@captype} is defined but there is no float, such as a
% minipage that fakes one for \cs{caption}, it, or a side caption, marks
% some unrelated box register instead.  That does no harm: a mark only
% ever reserves margin, and it is cleared when the next float starts.)
%
% With \cs{raggedbottom} the bottom floats sit right below the text rather
% than at the bottom of the page.  The notes may then also go beside the
% empty space below the floats, which is harmless; and as before, notes may
% already sit beside a bottom float that came up beside the margin of the
% text, whatever the mode.
%
% \section{Interaction with other packages}
% \subsection{memoir}
% There are no known issues with \Lopt{memoir} at present, provided that
% \cs{sidebar} is not used.
%
% \subsection{mparhack}
% \Lopt{mparhack} was designed to deal with the problem of margin notes
% showing up in the wrong margin because the left/right was decided before
% it was known exactly which page the note would be on.  Because we defer
% this decision to shipout time in this package, we are not susceptible
% to this problem, so \Lopt{mparhack} is no longer needed and should not
% be included (though I'm unaware whether it causes any actual problems).
%
% \subsection{Multiple columns}
% There is currently no support for multiple columns.
%
% \section{Coming attractions and known issues}
% Here is a list of things to possibly look forward to in a future version.
% If any of them are particularly important, please let me know.
% \begin{itemize}
% \item Use of pdf\TeX's \cs{pdfsavepos} and \cs{pdflastypos} for more
%   accurate margin placement.
% \item \cs{vadjust} to correct inconsistencies with \cs{@pageht}.
% \item Margin note placement is irrespective of vertical stretch.
%   Previously we gobbled any vertical stretch, but now that we have
%   fixed that bug, there's the possibility of wrong alignment since
%   we don't know where the positions will ultimately end up.  This
%   may be fixed by \cs{pdfsavepos} as well.
% \item Better interaction with floats.
%   (We can set a default one way
%   or the other and then allow a macro to override it (presumably with a
%   CS defined in terms of the box name/meaning, so as not to get in the
%   way of \LaTeX's use of the insert registers).  We would then add or
%   not add phantoms in the right spots.  We'd also need to shift all the
%   callout points by the size of the top figures (unless we're using
%   \cs{pdfsavepos}).)
% \end{itemize}
%
% \StopEventually{}
%
% \makeatletter
% \part*{Implementation}
% \section{Initial Setup}
% Make the |@|-sign into a letter for use in macro names.
%    \begin{macrocode}
%<*package>
\makeatletter
%    \end{macrocode}
%
% \begin{macros}{\MFX@debug}
% We have some optionally-included code for debugging.  \cs{MFX@debug}
% prints a new line followed by ``|MFX: |'' and then the message.
% We'll also ask for more error context in the debug mode.
%    \begin{macrocode}
%<*debug>
\def\MFX@debug{\message{^^JMFX:}\message}
\errorcontextlines=20
\def\MFX@mac#1{\expandafter\MFX@@mac\meaning#1>>>}
\def\MFX@@mac#1->{<<<}
\def\MFX@htdp#1{\ht#1=\the\ht#1, \dp#1=\the\dp#1}
%</debug>
%    \end{macrocode}
% \end{macros}
%
% The reader might begin to note at this point a convention we
% adopt throughout this package.  While we strive to avoid introducing
% new names as much as possible (with clever usages of \cs{expandafter}),
% any new names we do introduce will be prefixed
% by |\MFX@|, |\Mfx@|, or |\mfx@|, depending on the type of name.
% The all-capitol |\MFX@| is used for fully-constant macros.  The
% initial-caps |\Mfx@| is used for control sequences that are
% technically constant, but that refer to things that change, such
% as counters, token lists, dimension registers, etc.  Finally,
% the lowercase |\mfx@| is used for control sequences whose meaning
% changes dynamically (i.e. variable macros).
%
% \section{Options}
% Here we define the various package options.
% \iffalse - CURRENTLY UNIMPLEMENTED
% \begin{macros}{\ifmfx@ypos}
% The \Lopt{ypos} option signifies that we should use the pdf{\TeX}
% primitives \cs{pdfsavepos} and \cs{pdflastypos} to improve positioning
% of margin notes relative to their callouts.  This requires two
% passes to work, and the first time through, the margin notes
% will be positioned very na\"ively.
%    \begin{macrocode}
\newif\ifmfx@ypos
\DeclareOption{ypos}{\mfx@ypostrue}
%    \end{macrocode}
% \end{macros}
% END IFFALSE - \else There are no options yet. \fi
%
% Now we actually process the options.
%    \begin{macrocode}
\ProcessOptions\relax
%    \end{macrocode}
%
% \section{Variables}
% \begin{macros}{\mfx@marginlist}
% We need a place to store our list of marginal material.  We
% store material in this variable using insert registers and
% a variety of macros, to be explained later.
%    \begin{macrocode}
\let\mfx@marginlist\@empty
%    \end{macrocode}
% \end{macros}
%
% \begin{macros}{\mfx@inject,\Mfx@inject@insert}
% These are used to hijack \cs{marginpar} to inject arbitrary code
% into the output routine, rather than actually set a note.  To
% inject code, we unshift two copies of this dummy insert onto
% \cs{@freelist} for \cs{marginpar} to pull off.  We append
% whatever code we want to inject into \cs{mfx@inject} and then
% call \cs{marginpar}.  Then our custom \cs{@addmarginpar} will
% recognize the dummy inserts and run the code instead of setting
% a margin note.
%    \begin{macrocode}
\let\mfx@injected\@empty
\newinsert\Mfx@inject@insert
%    \end{macrocode}
% \end{macros}
%
% \begin{macros}{\Mfx@marginbox}
% While we're building the margin, we need to put it in a box
% before we can attach it to the main columm.
%    \begin{macrocode}
\newbox\Mfx@marginbox
%    \end{macrocode}
% \end{macros}
%
% \begin{macros}{\Mfx@marginpos@min,\Mfx@marginpos@max,\Mfx@marginspace}
% While we build up the margin piece boxes, we need to keep track of the
% possible range of positions.  The pair \cs{Mfx@marginpos@min} and
% \cs{Mfx@marginpos@max} are used to accumulate how much material has
% been added so far, with the difference that \cs{Mfx@marginpos@min} doesn't
% take into account compressible space, while \cs{Mfx@marginpos@max} does.
% Finally, \cs{Mfx@marginspace} is the amount of (incompressible) space
% since the last note, which allows skips to span margin phantoms.
%    \begin{macrocode}
\newdimen\Mfx@marginpos@min
\newdimen\Mfx@marginpos@max
\newdimen\Mfx@marginspace
%    \end{macrocode}
% \end{macros}
%
% \begin{macros}{\Mfx@marginheight}
% Because the margin height can be altered by, \cs{extendmargin},
% we must maintain a dimension for the height of the current margin.
% This dimension is reused in several different ways in the shipout-time
% margin building routines, keeping track of how much much space is left
% in (for the global passes) and the end position of the end of (for the
% piecewise passes) the current piece.
%    \begin{macrocode}
\newdimen\Mfx@marginheight
%    \end{macrocode}
% \end{macros}
%
% \begin{macros}{\mfx@marginstart,\mfx@marginpieces,
%     \Mfx@piece@content,\Mfx@piece@count,\ifmfx@in@phantom}
% These control sequences keep track of the margin phantoms.  When the
% margin is unblocked then \cs{mfx@marginstart} is not \cs{relax}.  When
% a phantom begins, the current \cs{mfx@marginstart} and page position
% are stored as a pair in \cs{mfx@marginpieces}, which is iterated over
% while building individual pieces of the margin.  We also define a box
% to keep the content of each margin piece, and a counter to keep track
% of how many pieces we have.  Finally, we define a switch for use in
% the second pass to indicate that we're inside a phantom.
%    \begin{macrocode}
\def\mfx@marginstart{0pt}
\let\mfx@marginpieces\@empty
\newbox\Mfx@piece@content
\newcount\Mfx@piece@count
\newif\ifmfx@in@phantom
%    \end{macrocode}
% \end{macros}
%
% \begin{macros}{\Mfx@mparshift}
% We store the current shift in a dimension register.
%    \begin{macrocode}
\newdimen\Mfx@mparshift
%    \end{macrocode}
% \end{macros}
%
% \section{User-configurable dimensions}
% We export a few dimensions that the user can redefine to tweak behavior.
% \begin{macros}{\marginheightadjustment}
% This length will be added to the total margin height of each page
% (the default is zero).
%    \begin{macrocode}
\newdimen\marginheightadjustment
%    \end{macrocode}
% \end{macros}
%
% \begin{macros}{\marginposadjustment}
% We will offset each margin note from its callout location by this length
% (the default is zero).
%    \begin{macrocode}
\newdimen\marginposadjustment
%    \end{macrocode}
% \end{macros}
%
% \begin{macros}{\marginsplitslack}
% When a deferrable note is split (mode \texttt{split}), we prefer a
% paragraph break that leaves at most this much of the room unused (but
% never more than half of it) to a line break.  The default is four lines of
% the normal text.  The value that counts is the one in force when a page is
% built, so set it outside groups, or globally.
%    \begin{macrocode}
\newdimen\marginsplitslack
\marginsplitslack=4\baselineskip
%    \end{macrocode}
% \end{macros}
%
% \section{Plan of attack}
% \subsection{\cs{marginpar}}
% The default sequence of events for a \cs{marginpar} is roughly the following
% (assuming no errors):
% \begin{verbatim}
% \marginpar:
%   let \@floatpenalty := (horizontal ? -10002 : -10003)
%   allocate inserts \@currbox and \@marbox from \@freelist
%   let \count\@marbox := -1 % signifies marginpar (not float)
%   if optional argument then \@xmpar else \@ympar
% \@xmpar:
%   \@savemarbox \@currbox := required argument
%   \@savemarbox \@marbox := optional argument
%   \@xympar
% \@ympar:
%   \@savemarbox \@currbox := required argument
%   copy \@marbox := \@currbox
%   \@xympar
% \@xympar:
%   append \@marbox to \@currlist
%   \end@float
% \end@float:
%   append \@currbox to \@currlist
%   if horizontal then following two lines are in \vadjust:
%     \penalty -10004
%     \penalty \@floatpenalty
% \end{verbatim}
%
% To get the rest of the picture, we need to peek into the
% output routine.  The pertinent parts are as follows (in
% vanilla \LaTeX):
% \begin{verbatim}
% \output:
%   if \outputpenalty < -10000 then
%     \@specialoutput
%   else
%     do regular output...
%   details for dealing with footnotes...
% \@specialoutput:
%   switch \outputpenalty:
%    case -10001: \@doclearpage
%    case -10004: set box \@holdpg := \vbox{\unvbox255}
%    case -10002 or -10003:
%     set box \@holdpg := \vbox{\unvbox\@holdpg \unvbox255}
%     let \@pageht := \ht\@holdpg, \@pagedp := \dp\@holdpg
%     \unvbox\@holdpg
%     pop \@currbox off of \@currlist
%     \@addmarginpar (assuming \count\@currbox <= 0)
% \@addmarginpar:
%   pop \@marbox off of \@currlist
%   free \@currbox and \@marbox back to \@freelist
%   if left-hand margin then let \@marbox := \@currbox
%   let \@tempdima := \@mparbottom - \@pageht + \ht\@marbox
%   if \@tempdima < 0 then let \@tempdima := 0
%   let \@mparbottom := \@pageht + \@tempdima + \dp\@marbox + \marginparpush
%   decrement \@tempdima := \@tempdima - \ht\@marbox
%   prepend \vskip\@tempdima to \@marbox
%   let \ht\@marbox := \dp\@marbox := 0
%   \kern -\@pagedp, \nointerlineskip
%   set an \hbox to \columnwidth (zero height/depth):
%     attach \@marbox to correct margin
%   set a \vbox with height 0 and depth \@pagedp
% \end{verbatim}
%
% We see from here that \cs{@addmarginpar} is the place where {\LaTeX}
% does the work of calculating the current page position and where the
% next note should go, and then actually puts it there.  We will need
% to completely replace this routine, but can leave everything else
% as is.
%
% \subsection{\cs{output}}
% While \LaTeX's margin routines end with \cs{@addmarginpar}, we must
% dig even deeper to apply our patch, since we need to insert some
% code to run during the \emph{main} output routine that ships out
% each page.  Thus, we'll expand ``\texttt{do regular output...}''
% from the previous \cs{output} listing.
% \begin{verbatim}
% do regular output...:
%   \@makecol
%   do { \@opcol \@startcolumn } while @fcolmade
% \@makecol:
%   set box \@outputbox := box255 (plus any footnotes)
%   let \@freelist := \@freelist + \@midlist, \@midlist := \@empty
%   \@combinefloats
%   add \@texttop and \@textbottom to \@outputbox (default no-op)
% \@opcol:
%   \@outputpage (or \@outputdblcol in twocolumn mode)
%   let \@mparbottom := \@textfloatsheight := 0
%   \@floatplacement
% \@startcolumn:
%   try to make a float column from \@deferlist, setting @fcolmade
%   if !@fcolmade then add floats from \@deferlist to next column
% \@combinefloats:
%   aggregate \@toplist floats into a box and prepend to \@outputbox
%   aggregate \@botlist floats into a box and append to \@coutputbox
%   free inserts from \@toplist and \@botlist
% \@outputpage:
%   ship out the page
%   reset a bunch of stuff
%   let \@colht := \textheight (in \@outputpage)
% \end{verbatim}
%
% We've seen two main times when action occurs: callout time and
% shipout time.  We proceed chronologically with our patches.
%
% \section{Callout-time patches}
% \begin{macros}{\@addmarginpar}
% The first thing we must modify is that at callout time, we need to
% get the inserts into \cs{mfx@marginlist}.  This should happen in the
% output routine so that we can get ahold of the current page
% position.  Even if we have a better idea of the page position
% (e.g. from pdf\TeX), we still might as well do this in the OR.
% In addition to actually setting the margin note, we also use this
% routine to inject arbitrary code into the OR (see \cs{MFX@inject}).
%    \begin{macrocode}
\def\@addmarginpar{%
  \@next\@marbox\@currlist{}\MFX@AssertionError
%<debug>\MFX@debug{addmarginpar (running insert) \@marbox/ \@currbox at
%<debug>    \the\c@page:\the\@pageht, marginlist=\MFX@mac\mfx@marginlist}%
%<debug>\MFX@debug{addmarginpar outputpenalty=\the\outputpenalty}%
  \MFX@getypos
  \expandafter\ifx\@marbox\Mfx@inject@insert
    \mfx@injected\global\let\mfx@injected\@empty
  \else
    \global\advance\mfx@placeid\@ne
    \mfx@note@locate
    \mfx@maybe@defer
  \fi
%<debug>\MFX@debug{addmarginpar (exit): marginlist=\MFX@mac\mfx@marginlist}%
}
%    \end{macrocode}
% \end{macros}
%
% \begin{macros}{\mfx@maybe@defer,\mfx@place@now,\mfx@defer@note}
% This is where ``keep notes on their anchor's page'' hooks in.  Every
% real marginpar is given a sequential id (\cs{mfx@placeid}) matching the
% id assigned at callout time (\cs{mfx@callid}); see \cs{mfx@anchorpage}.
% If a previous run recorded the page of this note's anchor (its callout)
% and that page lies \emph{after} the page we are currently shipping, then
% \cs{marginpar} drifted backwards across a page break: instead of setting
% the note here we hold it in \cs{mfx@deferred} until its anchor page comes
% around (see \cs{mfx@release@deferred}).  With no recorded anchor (first
% run, or a brand-new note) we behave exactly as before.
%    \begin{macrocode}
\def\mfx@maybe@defer{%
  \@ifundefined{mfx@ap@\the\mfx@placeid}%
    {\mfx@place@now}%
    {\ifnum\@nameuse{mfx@ap@\the\mfx@placeid}>\c@page
       \mfx@defer@note
     \else
       \mfx@place@now
     \fi}%
}
\def\mfx@place@now{%
  \MFX@cons\mfx@marginlist{%
    \noexpand\mfx@build@note\@currbox\@marbox{\mfx@ypos}%
    \noexpand\mfx@build@skip{\the\marginparpush}%
  }%
}
\def\mfx@defer@note{%
  \MFX@cons\mfx@deferred{%
    \noexpand\mfx@deferred@note\@currbox\@marbox
      {\@nameuse{mfx@ap@\the\mfx@placeid}}%
  }%
}
%    \end{macrocode}
% \end{macros}
%
% \begin{macros}{\MFX@cons,\MFX@snoc}
% In passing we'll define the cons macro, which fully-expands
% its second argument, but makes sure to only expand the first
% one once, so that any fragile control sequences in it are
% correctly protected.  We also define snoc, which prepends.
% Note that we could put the \cs{temp@} definition into a group
% if it was really gonna matter\ldots
%    \begin{macrocode}
\def\MFX@cons#1#2{%
  \edef\temp@{#2}%
  \expandafter\expandafter\expandafter\gdef
  \expandafter\expandafter\expandafter#1%
  \expandafter\expandafter\expandafter{\expandafter#1\temp@}%
}

\def\MFX@snoc#1#2{%
  \edef\temp@{#2}%
  \expandafter\expandafter\expandafter\gdef
  \expandafter\expandafter\expandafter#1%
  \expandafter\expandafter\expandafter{\expandafter\temp@#1}%
}
%    \end{macrocode}
% \end{macros}
%
% \begin{macros}{\MFX@run@clear}
% Finally, \cs{MFX@run@clear} is a quick trick to expand the
% contents of a macro and then clear it (to \cs{@empty}) before
% any of its tokens are consumed.
%     \begin{macrocode}
\def\MFX@run@clear#1{%
  \expandafter\global\expandafter\let\expandafter#1\expandafter\@empty#1%
}
%    \end{macrocode}
% \end{macros}
%
% \begin{macros}{\MFX@inject}
% As mentioned earlier, \cs{@addmarginpar} is also a hook for injecting
% arbitrary code into the output routine, i.e. to get a vertical position
% for blocking the margin.  We define \cs{MFX@inject} to facilitate this.
%    \begin{macrocode}
\def\MFX@inject#1{
  \expandafter\def\expandafter\@freelist\expandafter{%
    \expandafter\@elt\expandafter\Mfx@inject@insert
    \expandafter\@elt\expandafter\Mfx@inject@insert
    \@freelist}%
  \expandafter\def\expandafter\mfx@injected\expandafter{\mfx@injected#1}%
  \mfx@in@injecttrue
  \marginpar{}%
  \mfx@in@injectfalse
}
%    \end{macrocode}
% \end{macros}
%
% \begin{macros}{\MFX@getypos,\mfx@ypos}
% We now need to settle on a way to determine the vertical position.
% Someday this may be an option, and will depend on a variety
% of factors.  But for starters, we define the simplest version.
% Note the subtraction of \cs{Mfx@strutheight}.  Ideally we would simply
% grab a copy of \cs{@holdpg} from the middle of \cs{@specialoutput}
% and then discard the last box to figure out what height we're really
% at, since \cs{@holdpg} includes the box from the line we're currently
% on, and we want to be level with the \emph{top} of that box, rather
% than the baseline.  But since \cs{@holdpg} is accessible only deep
% within \cs{@specialoutput}, and it's not worth the risky job of
% performing surgery on it (which is unfortunately brittle if anyone
% else has a similar idea), we instead resort to this approximation.
% And since this should ultimately be only a fallback for when
% \cs{pdflastypos} isn't available, it's good enough.
% (NOTE: we might be able to use a \cs{vadjust} instead here?)
%    \begin{macrocode}
\def\MFX@getypos{%
  \dimen@\dimexpr\@pageht+\@pagedp+\marginposadjustment+\Mfx@mparshift\relax
  \ifnum\outputpenalty=-10002\relax
    \advance\dimen@-\Mfx@strutheight
  \fi
  \edef\mfx@ypos{\the\dimen@}%
  \global\Mfx@mparshift\z@
}
%    \end{macrocode}
% \end{macros}
%
% \begin{macros}{\marginpar,\Mfx@strutheight}
% We need to make sure \cs{Mfx@strutheight} gets defined somewhere,
% and the best time is probably right before the \cs{marginpar} does
% its work, since that will most likely ensure we're using the right
% font for the line.
%    \begin{macrocode}
\newdimen\Mfx@strutheight
\edef\marginpar{%
  \unexpanded{\setbox\@tempboxa\hbox{\strut}\Mfx@strutheight\ht\@tempboxa}%
  \unexpanded{\mfx@marginpar@callout}%
  \expandafter\unexpanded\expandafter{\marginpar}%
}
%    \end{macrocode}
% \end{macros}
%
% \section{Keeping notes on their anchor's page}
% A \cs{marginpar} is a float, and \LaTeX{} can attach it to the page it is
% busy assembling even when the page eventually breaks \emph{before} the line
% that called it---so a note whose callout sits at the top of a page can be
% set in the margin of the \emph{previous} page.  marginfix arranges notes
% within whatever page \LaTeX{} hands them to it, and can push overflow
% forwards, but on its own it has no notion of which page a note ``belongs''
% to, so it cannot repair this backwards drift.
%
% We fix it across two passes.  Each real marginpar is numbered in document
% order, and at callout time we record the page of that callout in the
% \texttt{.aux} file with a deferred \cs{write} (the same trick \cs{label}
% uses, so the page is captured at shipout).  On the next run that recorded
% \emph{anchor page} lets \cs{@addmarginpar} recognise a note that would be
% set too early and hold it back until its anchor page is shipped.  Because
% marginfix attaches margin material only after the body column is built,
% margin notes never move the body's page breaks, so anchor pages are stable
% and the process converges---usually in a single extra pass, requested
% through the kernel's own rerun machinery.
%
% \begin{macros}{\mfx@callid,\mfx@placeid,\ifmfx@in@inject,\mfx@deferred}
% Two counters number the marginpars: \cs{mfx@callid} at callout time and
% \cs{mfx@placeid} in the output routine.  They stay in lock-step because
% every real marginpar passes through each exactly once, in the same order.
% \cs{ifmfx@in@inject} is true only while \cs{MFX@inject} fakes a
% \cs{marginpar} for its own bookkeeping, so those are skipped.
% \cs{mfx@deferred} holds notes waiting for their anchor page.
%    \begin{macrocode}
\newcount\mfx@callid
\newcount\mfx@placeid
\newif\ifmfx@in@inject
\let\mfx@deferred\@empty
%    \end{macrocode}
% \end{macros}
%
% \begin{macros}{\mfx@cpage}
% \cs{mfx@cpage} expands to the raw (arabic) page number.  When it is
% \cs{let} to \cs{relax} inside \cs{protected@write} it survives the write's
% \cs{edef} unexpanded and then expands at shipout, giving us the page the
% callout actually landed on (regardless of any fancy \cs{thepage} format).
%    \begin{macrocode}
\def\mfx@cpage{\number\c@page}
%    \end{macrocode}
% \end{macros}
%
% \begin{macros}{\mfx@marginpar@callout}
% Prepended to \cs{marginpar}, this numbers each real callout and records its
% anchor page for the next run.  It also notes where in the source the
% callout is and that a new note is about to be set; both are for deferrable
% notes (section~\ref{sec:deferrable}).
%    \begin{macrocode}
\def\mfx@marginpar@callout{%
  \ifmfx@in@inject\else
    \global\advance\mfx@callid\@ne
    \global\mfx@fresh@notetrue
    \expandafter\xdef\csname mfx@loc@\the\mfx@callid\endcsname{\mfx@where}%
    \if@filesw
      \protected@write\@auxout{\let\mfx@cpage\relax}{%
        \string\mfx@anchorpage{\the\mfx@callid}{\mfx@cpage}}%
    \fi
  \fi
}
%    \end{macrocode}
% \end{macros}
%
% \begin{macros}{\mfx@anchorpage}
% This is what the \texttt{.aux} file calls back.  Normally (while the file
% is read at \cs{begin}\texttt{\{document\}}) it just stores the anchor page
% keyed by id.  During the kernel's end-of-run consistency check the aux is
% re-read with \cs{@newl@bel} \cs{let} to \cs{@testdef}; we detect that and,
% if an anchor page is new or has changed, set \cs{@tempswa} so the kernel
% issues its standard ``Label(s) may have changed.  Rerun'' warning---which
% \texttt{latexmk} and friends already act on.
%    \begin{macrocode}
\def\mfx@anchorpage#1#2{%
  \ifx\@newl@bel\@testdef
    \@ifundefined{mfx@ap@#1}%
      {\@tempswatrue}%
      {\def\reserved@a{#2}%
       \expandafter\ifx\csname mfx@ap@#1\endcsname\reserved@a\else
         \@tempswatrue
       \fi}%
  \else
    \global\@namedef{mfx@ap@#1}{#2}%
  \fi
}
%    \end{macrocode}
% \end{macros}
%
% \begin{macros}{\mfx@release@deferred,\mfx@release@one,\mfx@deferred@note}
% Run at the start of every \cs{MFX@combinefloats@before}.  We walk the
% deferred list: a note whose anchor page is still ahead of the page being
% shipped stays deferred; one whose page has arrived is released by prepending
% it to \cs{mfx@marginlist}, ahead of that page's own notes and at the top of
% the text block (its callout drifted from near the page top), so marginfix's
% normal building and stacking take over from there.  Released notes are
% collected in order and prepended as a block so several notes landing on the
% same page keep their document order.
%    \begin{macrocode}
\def\mfx@release@deferred{%
  \ifx\mfx@deferred\@empty\else
    \let\mfx@deferred@kept\@empty
    \let\mfx@released\@empty
    \let\mfx@deferred@note\mfx@release@one
    \mfx@deferred
    \global\let\mfx@deferred\mfx@deferred@kept
    \ifx\mfx@released\@empty\else
      \xdef\mfx@marginlist{%
        \unexpanded\expandafter{\mfx@released}%
        \unexpanded\expandafter{\mfx@marginlist}}%
    \fi
  \fi
}
\def\mfx@release@one#1#2#3{%
  \ifnum#3>\c@page
    \MFX@cons\mfx@deferred@kept{\noexpand\mfx@deferred@note#1#2{#3}}%
  \else
    \MFX@cons\mfx@released{%
      \noexpand\mfx@build@note#1#2{0pt}%
      \noexpand\mfx@build@skip{\the\marginparpush}}%
  \fi
}
%    \end{macrocode}
% \end{macros}
%
% \begin{macros}{\mfx@flush@deferred,\mfx@flush@one}
% A safety net: if a recorded anchor page is stale (e.g.\ the document shrank
% since the run that wrote it) a deferred note's page might never come around.
% \cs{dumpmargins} calls this at end of document to release everything that is
% still waiting, so a note is never silently dropped---worst case it lands a
% little late, exactly as an over-long note would today.
%    \begin{macrocode}
\def\mfx@flush@deferred{%
  \ifx\mfx@deferred\@empty\else
    \let\mfx@released\@empty
    \let\mfx@deferred@note\mfx@flush@one
    \mfx@deferred
    \global\let\mfx@deferred\@empty
    \ifx\mfx@released\@empty\else
      \xdef\mfx@marginlist{%
        \unexpanded\expandafter{\mfx@released}%
        \unexpanded\expandafter{\mfx@marginlist}}%
    \fi
  \fi
}
\def\mfx@flush@one#1#2#3{%
  \MFX@cons\mfx@released{%
    \noexpand\mfx@build@note#1#2{0pt}%
    \noexpand\mfx@build@skip{\the\marginparpush}}%
}
%    \end{macrocode}
% \end{macros}
%
% \section{Deferrable notes}
% \label{sec:deferrable}
% A note is deferrable if its text contains \cs{marginnotedeferrable}.  We
% need to know that in the output routine, where all we have of a note is
% its pair of box registers, so \cs{marginnotedeferrable} records it in a
% control sequence named after the register numbers,
% |\mfx@cls@|\meta{number}.  It runs while
% the note's text is being typeset into its box, in \cs{@savemarbox}, and
% then \cs{@currbox} and \cs{@marbox} name the two registers that the output
% routine will later see.  (Which of the two ends up first in
% \cs{mfx@marginlist} depends on the order \LaTeX{} pops them, so we mark
% both.)
%
% Registers are reused, so the mark of an earlier note must not stick to a
% later one.  We clear both marks when a new note starts to be typeset, which
% we recognise by \cs{@marginparreset}: \LaTeX{} calls it at the start of
% every \cs{@savemarbox}.  With the optional argument of \cs{marginpar} there
% are two calls for one note, so we clear only on the first, flagged by
% \cs{mfx@marginpar@callout}; a mark in either text then marks the note.
% The same hook tells \cs{marginnotedeferrable} that it is inside a margin
% note.
%
% \begin{macros}{\ifmfx@fresh@note,\ifmfx@in@note,\mfx@note@reset,
%     \MFX@forget}
% \cs{MFX@forget} also clears the flag that the note was moved whole from
% the previous page, which we need for placing it on the next
% (\cs{MFX@carrylate}).  These are only settings of flags, so a document
% without deferrable notes is typeset exactly as before.
%    \begin{macrocode}
\newif\ifmfx@fresh@note
\newif\ifmfx@in@note
\g@addto@macro\@marginparreset{\mfx@note@reset}
\def\mfx@note@reset{%
  \mfx@in@notetrue
  \ifmfx@fresh@note
    \global\mfx@fresh@notefalse
    \MFX@forget\@currbox
    \MFX@forget\@marbox
  \fi
}
\def\MFX@forget#1{%
  \expandafter\global\expandafter\let
    \csname mfx@cls@\number#1\endcsname\relax
  \expandafter\global\expandafter\let
    \csname mfx@carried@\number#1\endcsname\relax
}
%    \end{macrocode}
% \end{macros}
%
% \begin{macros}{\marginnotedeferrable,\MFX@markdeferrable,\MFX@notefont,
%     \MFX@islow}
% The mark itself, and the test for it: \cs{MFX@islow}\marg{box} sets
% \cs{@tempswa} to whether the note in \emph{box} is deferrable.  The mark
% also records the font and \cs{baselineskip} at that point, in
% |\mfx@nf@|\meta{number} and |\mfx@nb@|\meta{number}, for the cues of a
% split note, which are typeset in the output routine, where the note's
% font is long gone.  We record the font by its NFSS attributes rather
% than by \cs{font}, so that font commands in the cue text work relative to
% it; and the \cs{baselineskip} as it is, since \cs{baselinestretch} in the
% output routine may differ from the note's.  \cs{MFX@notefont}\marg{box}
% selects them again, if there are any.
%    \begin{macrocode}
\def\MFX@markdeferrable#1{%
  \expandafter\gdef\csname mfx@cls@\number#1\endcsname{}%
  \expandafter\xdef\csname mfx@nf@\number#1\endcsname{%
    \noexpand\fontsize{\f@size}{\the\baselineskip}%
    \noexpand\usefont{\f@encoding}{\f@family}{\f@series}{\f@shape}}%
  \expandafter\xdef\csname mfx@nb@\number#1\endcsname{\the\baselineskip}%
}
\def\MFX@notefont#1{%
  \expandafter\ifx\csname mfx@nf@\number#1\endcsname\relax\else
    \csname mfx@nf@\number#1\endcsname
    \baselineskip\csname mfx@nb@\number#1\endcsname\relax
  \fi
}
\def\marginnotedeferrable{%
  \ifmfx@in@note
    \MFX@markdeferrable\@currbox
    \MFX@markdeferrable\@marbox
  \else
    \PackageWarning{marginfix}{\string\marginnotedeferrable\space outside a
      margin note ignored}%
  \fi
}
\def\MFX@islow#1{%
  \expandafter\ifx\csname mfx@cls@\number#1\endcsname\relax
    \@tempswafalse
  \else
    \@tempswatrue
  \fi
}
%    \end{macrocode}
% \end{macros}
%
% \begin{macros}{\mfx@where,\mfx@note@locate,\MFX@loc}
% The warnings name notes by where they are in the source.  At callout time
% \cs{mfx@marginpar@callout} stores the file and line under the callout's
% id, and \cs{@addmarginpar} moves it to the note's register once it knows
% the register, the same way the anchor page is matched to the note.  The
% line is the one where the \cs{marginpar} is expanded, that is, where the
% note's text ends.  In the main file \cs{CurrentFile} is empty, so we use
% the job name instead.  \cs{MFX@loc}\marg{box} expands to the location, or
% to a question mark if we don't know it.
%    \begin{macrocode}
\def\mfx@where{%
  \ifx\CurrentFile\@empty
    \jobname
  \else\ifx\CurrentFile\@undefined
    \jobname
  \else
    \ifx\CurrentFilePath\@empty\else\CurrentFilePath/\fi\CurrentFile
  \fi\fi
  :\the\inputlineno}
\def\mfx@note@locate{%
  \expandafter\global\expandafter\let
    \csname mfx@boxloc@\number\@currbox\expandafter\endcsname
    \csname mfx@loc@\the\mfx@placeid\endcsname
  \expandafter\global\expandafter\let
    \csname mfx@loc@\the\mfx@placeid\endcsname\relax
}
\def\MFX@loc#1{%
  \expandafter\ifx\csname mfx@boxloc@\number#1\endcsname\relax
    ?%
  \else
    \csname mfx@boxloc@\number#1\endcsname
  \fi
}
%    \end{macrocode}
% \end{macros}
%
% \section{Shipout-time patches}
% \begin{macros}{\@combinefloats}
% We need to patch in somewhere before \cs{@combinefloats} at the latest,
% so that any heights calculated from \cs{@pageht} are correct---otherwise
% the top figures will confuse us.  So we'll start by simply adding our
% own \cs{MFX@combinefloats@before} at the very beginning of \cs{@combinefloats}
%    \begin{macrocode}
\expandafter\def\expandafter\@combinefloats\expandafter{\expandafter
  \MFX@combinefloats@before\@combinefloats}
%    \end{macrocode}
% Since the 2020-10-01 release, the kernel's \cs{@makecol} no longer calls
% \cs{@combinefloats} directly but through \cs{@outputbox@attachfloats},
% which was \cs{let} to \cs{@combinefloats} before we patched it.  The
% kernel's first-aid file repointed it for v1.2 of this package only, so we
% must do it ourselves; otherwise our hook never runs under the standard
% classes (\cs{memoir} has its own \cs{@makecol} that still calls
% \cs{@combinefloats}), the notes are never attached, and the document ends
% with ``Too many unprocessed floats'' and ``lost some margin notes''.
%    \begin{macrocode}
\ifdefined\@outputbox@attachfloats
  \let\@outputbox@attachfloats\@combinefloats
\fi
%    \end{macrocode}
% \end{macros}
%
% \begin{macros}{\MFX@combinefloats@before}
% \cs{MFX@combinefloats@before} is then responsible for picking the
% needed notes from \cs{mfx@marginlist}, building them into a box, and
% attaching that box onto the correct side of \cs{@outputbox}.  We also
% add any global \cs{marginheightadjustment} to \cs{Mfx@marginheight}
% before building the margin, and then reset it back to zero at the end.
% This allows any calls to \cs{extendmargin} during the page itself to
% work as expected.
%    \begin{macrocode}
\def\MFX@combinefloats@before{%
  \advance\Mfx@marginheight\marginheightadjustment
  \mfx@release@deferred
  \MFX@buildmargin
  \MFX@attachmargin
  \global\Mfx@marginheight\z@
}
%    \end{macrocode}
% \end{macros}
%
% \begin{macros}{\MFX@attachmargin}
% We'll start with the second half of \cs{MFX@combinefloats@before},
% since it's simpler.  We need to do several things here.
%    \begin{macrocode}
\def\MFX@attachmargin{%
%<debug>\MFX@debug{attachmargin}%
%    \end{macrocode}
% We start by moving the reference point of \cs{Mfx@marginbox} to the top.
%    \begin{macrocode}
%<debug>\MFX@debug{attachmargin: \MFX@htdp\@outputbox, \MFX@htdp\Mfx@marginbox}%
  \setbox\Mfx@marginbox\vtop{%
    \vskip\z@\unvbox\Mfx@marginbox}%
%    \end{macrocode}
% Next we need to figure out which side of \cs{@outputbox} to attach
% the \cs{Mfx@marginbox} on.  We now use \cs{columnwidth} instead of
% \cs{wd}\cs{@outputbox} to set the right-hand margins, since tufte-\LaTeX
% sometimes makes too-wide output boxes.  If this becomes a problem,
% we'll need to consider making this configurable elsewhere.  We should
% also pay attention to whether adding a box at the top of \cs{@outputbox}
% might have unintended consequences w.r.t. any glue being retained that
% should have been swallowed.  This will require further investigation.
%    \begin{macrocode}
  \setbox\@outputbox\vbox{%
    \begingroup
    \setbox\@tempboxa\vbox{%
      \hbox{%
        \if\MFX@leftmargin
          \llap{\box\Mfx@marginbox\hskip\marginparsep}%
        \else
          \hskip\columnwidth
          \rlap{\hskip\marginparsep\box\Mfx@marginbox}%
        \fi
      }}%
    \ht\@tempboxa\z@
    \dp\@tempboxa\z@
    \box\@tempboxa
    \endgroup
    \unskip
    \unvbox\@outputbox
  }%
}
%    \end{macrocode}
% \end{macros}
%
% \begin{macros}{\MFX@buildmargin}
% When \cs{MFX@buildmargin} is called, we have a list of tokens in
% \cs{mfx@marginlist} that need to be processed: combinations of
% \cs{mfx@build@note}, \cs{mfx@build@skip}, and \cs{mfx@build@clear},
% with various parameters to indicate what material still needs to be
% set in the margin.  This macro therefore must pull off the first $n>0$
% of these commands to set on the current page ($n$ must be positive to
% prevent infinite loops), and leave the rest to be deferred.  Notes are
% taken in order, except that deferrable notes may give way to the ordinary
% notes after them, on an opt-in basis (section~\ref{sec:priority@down}).
% The typeset
% material will be left in \cs{Mfx@marginbox}, which must have the same
% height as \cs{@outputbox} (although because we \cs{unvbox}\cs{@outputbox}
% in \cs{MFX@attachmargin}, we can't guarantee that this will correctly
% line up the notes with their callouts).  This procedure happens in
% four passes.  But first, we initialize \cs{Mfx@marginheight} to
% \cs{@colroom}, which is the height of the page minus any floats that
% have been added to the top or bottom (these floats may extend into the
% margins: in the future we may look into detecting this and using the
% whole page, with overwide floats blocked off as phantoms).  We add
% \cs{@colroom} rather than assigning it because any global or per-page
% adjustments have already been added to \cs{Mfx@marginheight}.  We
% can then close out any still-open margin pieces (this is the typical
% case, where the margin is not blocked across a page boundary, so that
% \cs{mfx@marginstart} will hold a position, rather than \cs{relax}).
% After this, \cs{Mfx@marginheight} is no longer necessary, so we reuse
% it for keeping track of available space in individual margin pieces.
% With \cs{marginnotebesidefloats}\texttt{\{use\}}, the last piece also
% reaches down beside the bottom floats, if none of them extends into the
% margin (section~\ref{sec:floats@impl}).
%    \begin{macrocode}
\def\MFX@buildmargin{%
  \advance\Mfx@marginheight\@colroom
  \csname MFX@beside@\mfx@besidefloats\endcsname
  \ifx\mfx@marginstart\relax
  \else
    \MFX@cons\mfx@marginpieces{%
      \noexpand\@elt{\mfx@marginstart}{\the\Mfx@marginheight}}%
    \gdef\mfx@marginstart{0pt}%
    \global\advance\Mfx@piece@count\@ne
  \fi
%<debug>\MFX@debug{buildmargin: marginheight=\the\Mfx@marginheight,
%<debug>           marginlist=\MFX@mac\mfx@marginlist,
%<debug>           marginpieces=\MFX@mac\mfx@marginpieces}%
%    \end{macrocode}
% We now execute the four passes.  First is a global downward pass,
% whose purpose is to determine the maximum number of notes (and other
% material) that can fit in the margin, taking any phantoms into
% consideration.  Every note identified by \cs{MFX@buildmargin@down}
% is guaranteed to show up on this page, so we free its inserts back
% to \cs{@freelist}.  The second pass is the global upward pass, in
% which we determine the lowest possible margin piece each note may
% go into without causing lower notes to fall off the bottom.  The
% third pass is the piecewise downward pass.  For each piece, we
% figure out where in the piece each note will go by inserting
% compressible spaces between the notes.  If a note is called out
% past the end of the piece and does not need to go into the piece
% (as determined by pass 2), it will be deferred to a later piece.
% The fourth pass is the piecewise upward pass, in which the compressible
% spaces are shrunk just enough to fit everything into the piece.  The
% last two (piecewise) passes both occur in each piece before the next
% piece is addressed.  The whole process is bypassed if there are
% no eligible margin pieces.
%    \begin{macrocode}
  \ifx\mfx@marginpieces\@empty\else
    \MFX@buildmargin@down
    \MFX@buildmargin@up
    \MFX@buildmargin@pieces
  \fi
}
%    \end{macrocode}
% \end{macros}
%
% \subsection{Room beside bottom floats}
% \label{sec:floats@impl}
% The margin box hangs from the top of the text, before \cs{@combinefloats}
% adds the floats, and its last piece ends at \cs{@colroom}, the height of
% the text.  The bottom floats come after that: \cs{@cflb} appends
% \cs{textfloatsep}, then each float followed by \cs{floatsep}, and takes
% the last \cs{floatsep} back.  So to let the notes continue beside them we
% need only extend the last piece by that much, which is what
% \cs{MFX@beside@use} does; the later passes then place notes there like
% anywhere else.  If the margin is blocked at the end of the page
% (\cs{blockmargin} without \cs{unblockmargin}) there is no last piece to
% extend, and nothing changes.
%
% We only extend if no bottom float extends into the margin, and then
% beside all of them; a float that does reserves the whole area, as before.
% Giving each float its own piece would let notes go beside the narrow
% ones, but bottom floats are seldom more than one or two, and the simpler
% rule leaves less to go wrong.
%
% Whether a float extends into the margin we have to find out from its
% box, and the box itself says little: a float box is as wide as the
% column whatever is in it.  A table wider than the text is an overfull
% line in a box of the column's width, and the lines of
% \texttt{adjustwidth} with a negative margin are wider than the column
% but shifted, so the box's width, the largest width plus shift, may still
% be the column's.  So we look at the lines (\cs{MFX@floatwide}): a line
% wider than the column, or overfull, makes the float wide.  That catches
% overfull tables, and text set wider on either side, since a line that
% sticks out on the left is wider than the column too.  We can't use the
% shifts themselves: \cs{lastbox} drops them.
%
% We walk the float's list from the top with \cs{vsplit}, one piece at a
% time, and take each piece apart from the bottom with \cs{lastbox},
% \cs{unskip}, \cs{unkern} and \cs{unpenalty}, recognising the items with
% e-\TeX's \cs{lastnodetype}; boxes in the list (a minipage, the inner box
% of the float) we walk in turn, down to a few levels.  Taking the whole
% list apart from the bottom would be simpler, but \TeX{} can't take a
% \cs{write} or a colour change out of a list, and floats are full of
% them: the caption writes to the list of tables, a \cs{label} writes to
% the \texttt{.aux} file, and with colour the float's own box ends with a
% colour change.  A piece from \cs{vsplit} ends at the next place the list
% may break, so a \cs{write} with a break before it is a piece of its own,
% which we can leave where it is.  But the caption's \cs{write} comes right
% after the table or picture, and the colour change right after the inner
% box of the float, with no break between.  So in the mode \texttt{use}
% we put a \cs{penalty}|0| before them, as the float is built: before every
% \cs{protected@write} in vertical mode inside a float (that is how
% \cs{label}, \cs{addcontentsline} and \cs{index} write), and after the
% inner box of the float, in \cs{@endfloatbox}.  Inside a float means that
% \cs{@captype} is defined and \cs{@floatpenalty} is negative, as it is
% throughout the body of a float, so that a \cs{captionof} on the page, a
% figure in \pkg{beamer} or a table that breaks across pages get no
% penalties.  A penalty is a place to break, but changes no size (unlike
% glue or a kern, which would end the depth of the box before it), and a
% float is never broken, so the float looks the same.  Anything we still
% can't take apart, and that has a size (a rule, or a box right before
% some other \cs{special}), makes the float wide, to be safe; and so does
% a list nested too deep.  If another package has changed
% \cs{@endfloatbox}, we leave it alone, and with colour every float is
% then wide; likewise if a package changes \cs{protected@write} in the
% document, or \cs{nofiles} turns it into a bare \cs{write}.  So when a
% float is wide for such a reason rather than for a wide line, we say so
% in the log (\cs{PackageInfo}), which helps to find out why notes don't
% go beside it.  The cost is one copy of each bottom float per page, and
% only in the mode \texttt{use}.
%
% The lines of a float may hold boxes too, and we look at those at the end
% of a line, from the right (\cs{MFX@fw@hlist}); see below for which.
%
% What \TeX{} can't see at all is ink outside the boxes: a \cs{rlap} or
% \cs{llap}, a \cs{makebox} that lets its content stick out, negative
% space that pulls text out of its line, or a drawing beyond its bounding
% box.  Such a float has to say \cs{marginnotereservefloat}.  \pkg{memoir}'s
% side captions are of that kind, so we mark the floats that use one,
% through the environment hooks of the \LaTeX{} kernel.  We don't try to
% tell which margin a side caption is in: reserving too much is only a
% waste.
%
% With \cs{raggedbottom} the text is not stretched to \cs{@colroom}, and the
% bottom floats follow the text directly.  The extended piece then also
% covers the empty space below them, which does no harm, and a narrow float
% may come up beside the margin of the text, which was already so.
%
% \begin{macros}{\MFX@beside@reserve,\MFX@beside@use,\MFX@bottomroom}
% The mode \texttt{reserve} does nothing, so the default layout is as
% before.  \cs{MFX@bottomroom} leaves in \cs{dimen@} the room beside the
% floats in \cs{@botlist}, or zero if one of them extends into the margin.
%    \begin{macrocode}
\def\MFX@beside@reserve{}
\def\MFX@fw@line{line}
\def\MFX@beside@use{%
  \ifx\@botlist\@empty\else
    \MFX@bottomroom
    \advance\Mfx@marginheight\dimen@
  \fi
}
\def\MFX@bottomroom{%
  \dimen@\z@
  \count@\z@
  \@tempswafalse
  \def\@elt##1{%
    \advance\count@\@ne
    \advance\dimen@\dimexpr\ht##1+\dp##1\relax
    \expandafter\ifx\csname mfx@wide@\number##1\endcsname\relax
      \if@tempswa\else
        \MFX@floatwide##1%
        \ifmfx@fw@wide
          \@tempswatrue
          \ifx\mfx@fw@why\MFX@fw@line\else
            \PackageInfo{marginfix}{No notes beside the bottom floats of
              page \thepage:\MessageBreak one of them
              holds \mfx@fw@why.\MessageBreak If it doesn't, a package may
              have changed\MessageBreak \string\protected@write\space or
              \string\@endfloatbox\@gobble}%
          \fi
        \fi
      \fi
    \else
      \@tempswatrue
    \fi}%
  \@botlist
  \let\@elt\relax
  \advance\dimen@\textfloatsep
  \advance\dimen@\numexpr\count@-\@ne\relax\floatsep
  \if@tempswa
    \dimen@\z@
  \fi
}
%    \end{macrocode}
% \end{macros}
%
% \begin{macros}{\MFX@floatwide,\ifmfx@fw@wide,\MFX@fw@list,\MFX@fw@peel,
%     \MFX@fw@items,\MFX@fw@hbox,\MFX@fw@vbox}
% \cs{MFX@floatwide}\marg{box} sets \cs{ifmfx@fw@wide} (globally) to
% whether the float in \emph{box} extends into the margin, as described
% above.  It works in a group, on a copy, with the warnings about bad boxes
% turned off, since we repack lines and split lists to find out.  It
% compares widths with the width of a column as \LaTeX\ sets it
% (\cs{mfx@fw@colwidth}), not with \cs{columnwidth}: we are in the output
% routine, inside whatever groups were open where the page broke, and
% one of them may have set a wider \cs{columnwidth}, as an environment
% that widens the text into the margin may do.
% \cs{MFX@fw@list}\marg{box} walks the list in \emph{box} a level deeper,
% in a group of its own, so that each level has its own registers.
%    \begin{macrocode}
\newif\ifmfx@fw@wide
\newcount\Mfx@fw@depth
\newcount\Mfx@fw@pieces
\newbox\Mfx@fw@rest
\newbox\Mfx@fw@piece
\newbox\Mfx@fw@item
\def\MFX@fw@maxdepth{8}
\def\MFX@floatwide#1{%
  \begingroup
    \global\mfx@fw@widefalse
    \hbadness\@M \hfuzz\maxdimen
    \vbadness\@M \vfuzz\maxdimen
    \splittopskip\z@skip
    \splitmaxdepth\maxdimen
    \Mfx@fw@depth\z@
    \Mfx@fw@pieces\z@
    \gdef\mfx@fw@why{}%
    \dimen@\textwidth
    \if@twocolumn
      \advance\dimen@-\columnsep
      \divide\dimen@\tw@
    \fi
    \edef\mfx@fw@colwidth{\the\dimen@}%
    \setbox\Mfx@fw@item\copy#1%
    \MFX@fw@vbox
  \endgroup
}
\def\MFX@fw@list{%
  \begingroup
    \advance\Mfx@fw@depth\@ne
    \ifnum\Mfx@fw@depth>\MFX@fw@maxdepth\relax
      \MFX@fw@setwide{boxes nested too deep}%
    \else
      \setbox\Mfx@fw@rest\vbox{\unvbox\Mfx@fw@item}%
      \MFX@fw@peel
    \fi
  \endgroup
}
%    \end{macrocode}
% Peeling: \cs{vsplit} to zero height takes the list up to its first
% legal break.  What is left in a piece after we have taken apart all we
% can must have no size.  We also stop after very many pieces, to be safe.
%    \begin{macrocode}
\def\MFX@fw@peel{%
  \let\mfx@fw@next\relax
  \ifvoid\Mfx@fw@rest\else
    \ifmfx@fw@wide\else
      \global\advance\Mfx@fw@pieces\@ne
      \ifnum\Mfx@fw@pieces>10000
        \MFX@fw@setwide{a list too long to walk}%
      \else
        \setbox\Mfx@fw@piece\vsplit\Mfx@fw@rest to\z@
        \setbox\Mfx@fw@piece\vbox{\unvbox\Mfx@fw@piece\MFX@fw@items}%
        \ifdim\dimexpr\ht\Mfx@fw@piece+\dp\Mfx@fw@piece\relax>\z@
          \MFX@fw@setwide{an item it can't take apart}%
        \fi
        \ifdim\wd\Mfx@fw@piece>\z@
          \MFX@fw@setwide{an item it can't take apart}%
        \fi
        \let\mfx@fw@next\MFX@fw@peel
      \fi
    \fi
  \fi
  \mfx@fw@next
}
%    \end{macrocode}
% Taking a piece apart from the bottom, inside the box being built.
%    \begin{macrocode}
\def\MFX@fw@items{%
  \let\mfx@fw@next\MFX@fw@items
  \ifcase\numexpr\lastnodetype+1\relax
    \let\mfx@fw@next\relax                    % empty
  \or\let\mfx@fw@next\relax                   % char
  \or\setbox\Mfx@fw@item\lastbox\MFX@fw@hbox  % hlist
  \or\setbox\Mfx@fw@item\lastbox\MFX@fw@vbox  % vlist
  \or\let\mfx@fw@next\relax\or\let\mfx@fw@next\relax % rule, ins
  \or\let\mfx@fw@next\relax\or\let\mfx@fw@next\relax % mark, adjust
  \or\let\mfx@fw@next\relax\or\let\mfx@fw@next\relax % ligature, disc
  \or\let\mfx@fw@next\relax\or\let\mfx@fw@next\relax % whatsit, math
  \or\unskip                                  % glue
  \or\unkern                                  % kern
  \or\unpenalty                               % penalty
  \else\let\mfx@fw@next\relax
  \fi
  \ifmfx@fw@wide\let\mfx@fw@next\relax\fi
  \mfx@fw@next
}
%    \end{macrocode}
% A line is too wide if it is wider than the column, or if its content
% doesn't fit in it, even with a point to spare: repacked to that width
% it is overfull, and then \TeX{} sets \cs{badness} to 1000000.  A box in
% the list is too wide if it is wider than the column; else we walk it.
%    \begin{macrocode}
\def\MFX@fw@hbox{%
  \ifdim\wd\Mfx@fw@item>\dimexpr\mfx@fw@colwidth+\p@\relax
    \MFX@fw@setwide{line}%
  \else
    \setbox\Mfx@fw@item
      \hbox to\dimexpr\wd\Mfx@fw@item+\p@\relax{\unhbox\Mfx@fw@item}%
    \ifnum\badness>\@M
      \MFX@fw@setwide{line}%
    \else
      \MFX@fw@hlist
    \fi
  \fi
}
\def\MFX@fw@vbox{%
  \ifdim\wd\Mfx@fw@item>\dimexpr\mfx@fw@colwidth+\p@\relax
    \MFX@fw@setwide{line}%
  \else
    \MFX@fw@list
  \fi
}
%    \end{macrocode}
% Within a line we look at the boxes at its end: we take the line apart
% from the right, glue, kerns, penalties and boxes, and look at each box
% as above, until we meet anything else, such as a character, a rule or a
% formula.  That is not a reason to call the float wide, since every line
% of text ends so; we simply stop.  This reaches a minipage, a
% \cs{parbox} or a table with the position \texttt{t} or \texttt{b} that
% ends its line, which is how they usually stand in a float; but not one
% centred on its line (the default), since \LaTeX{} puts that in a formula
% (\cs{vcenter}), nor anything followed by text, nor the inside of an
% \cs{fbox}, which ends with a rule.
%    \begin{macrocode}
\def\MFX@fw@hlist{%
  \begingroup
    \advance\Mfx@fw@depth\@ne
    \ifnum\Mfx@fw@depth>\MFX@fw@maxdepth\relax
      \MFX@fw@setwide{boxes nested too deep}%
    \else
      \setbox\Mfx@fw@item\hbox{\unhbox\Mfx@fw@item\MFX@fw@hitems}%
    \fi
  \endgroup
}
\def\MFX@fw@hitems{%
  \let\mfx@fw@next\relax
  \ifcase\numexpr\lastnodetype+1\relax
  \or                                                      % char
  \or\setbox\Mfx@fw@item\lastbox\MFX@fw@hbox
     \let\mfx@fw@next\MFX@fw@hitems                        % hlist
  \or\setbox\Mfx@fw@item\lastbox\MFX@fw@vbox
     \let\mfx@fw@next\MFX@fw@hitems                        % vlist
  \or\or\or\or\or\or\or\or                                % rule ... math
  \or\unskip\let\mfx@fw@next\MFX@fw@hitems                 % glue
  \or\unkern\let\mfx@fw@next\MFX@fw@hitems                 % kern
  \or\unpenalty\let\mfx@fw@next\MFX@fw@hitems              % penalty
  \fi
  \ifmfx@fw@wide\let\mfx@fw@next\relax\fi
  \mfx@fw@next
}
%    \end{macrocode}
% \cs{MFX@fw@setwide}\marg{reason} calls the float wide, keeping the first
% reason in \cs{mfx@fw@why}, for the message in \cs{MFX@bottomroom}.
%    \begin{macrocode}
\def\MFX@fw@setwide#1{%
  \ifmfx@fw@wide\else
    \global\mfx@fw@widetrue
    \gdef\mfx@fw@why{#1}%
  \fi
}
%    \end{macrocode}
% \end{macros}
%
% \begin{macros}{\MFX@fw@break,\MFX@fw@writebreak,\MFX@kernel@endfloatbox}
% The penalties that let us take a float apart, put in as it is built,
% only in the mode \texttt{use}.  We patch \cs{@endfloatbox} and
% \cs{protected@write} at the start of the document, after the other
% packages; \cs{@endfloatbox} only if it is still the kernel's.
%    \begin{macrocode}
\def\MFX@fw@usename{use}
\def\MFX@fw@break{%
  \ifx\mfx@besidefloats\MFX@fw@usename
    \penalty\z@
  \fi
}
\def\MFX@fw@writebreak{%
  \ifvmode
    \ifx\@captype\@undefined\else
      \ifnum\@floatpenalty<\z@
        \MFX@fw@break
      \fi
    \fi
  \fi
}
\def\MFX@kernel@endfloatbox{%
  \par\vskip\z@skip
  \@minipagefalse
  \outer@nobreak
  \egroup
  \color@endbox
}
\AtBeginDocument{%
  \ifx\@endfloatbox\MFX@kernel@endfloatbox
    \def\@endfloatbox{%
      \par\vskip\z@skip
      \@minipagefalse
      \outer@nobreak
      \egroup
      \MFX@fw@break
      \color@endbox}%
  \fi
  \let\MFX@orig@protected@write\protected@write
  \def\protected@write{\MFX@fw@writebreak\MFX@orig@protected@write}%
}
%    \end{macrocode}
% \end{macros}
%
% \begin{macros}{\marginnotereservefloat,\MFX@reservefloat}
% A float is marked in |\mfx@wide@|\meta{number}, named after its box
% register, as deferrable notes are.  Inside a float, \cs{@currbox} is that
% register, and \cs{@captype} is defined.  The mark is cleared when a float
% starts, in \cs{@floatboxreset}, since registers are reused.  The side
% captions mark their float silently, since they may also be used outside
% a float.
%    \begin{macrocode}
\g@addto@macro\@floatboxreset{%
  \expandafter\global\expandafter\let
    \csname mfx@wide@\number\@currbox\endcsname\relax}
\def\MFX@reservefloat{%
  \ifx\@captype\@undefined\else
    \expandafter\gdef\csname mfx@wide@\number\@currbox\endcsname{}%
  \fi
}
\def\marginnotereservefloat{%
  \ifx\@captype\@undefined
    \PackageWarning{marginfix}{\string\marginnotereservefloat\space
      outside a float ignored}%
  \else
    \MFX@reservefloat
  \fi
}
\ifdefined\AddToHook
  \AddToHook{env/sidecaption/begin}{\MFX@reservefloat}
  \AddToHook{env/sidecontcaption/begin}{\MFX@reservefloat}
  \AddToHook{env/sidelegend/begin}{\MFX@reservefloat}
  \AddToHook{env/sidenamedlegend/begin}{\MFX@reservefloat}
\fi
%    \end{macrocode}
% \end{macros}
%
% \subsection{First pass: global downward}
% \begin{macros}{\MFX@buildmargin@down,\mfx@pieceheights}
% The first step is the global ``down'' step, in which we move the
% notes that will fit on the current page into \cs{mfx@marginout} in
% reverse order (to prepare for the second, upward, pass), and anything
% that doesn't fit is deferred back into \cs{mfx@marginlist}.  We do
% this by changing the meaning of \cs{mfx@build@note}, \cs{mfx@build@skip},
% and \cs{mfx@build@clear}, which delimit the different types of material
% in \cs{mfx@marginlist}.  Note that as we continue processing, these
% macros will change from time to time (i.e. changing \cs{mfx@build@skip}
% to actually doing something once we find a note, rather than gobbling
% so as to remove skips at page boundaries; or changing them to save
% material back onto \cs{mfx@marginlist} once the margin fills up).
% The first thing we need to do is iterate over the piece positions
% to get the list of heights.
%    \begin{macrocode}
\def\MFX@buildmargin@down{%
  \let\mfx@pieceheights\@empty
  \def\@elt##1##2{%
    \MFX@cons\mfx@pieceheights{\noexpand\@elt{\the\dimexpr##2-##1}}}%
  \mfx@marginpieces
%    \end{macrocode}
% Then the priority mode (\cs{marginnotepriority}) gets its say.  It sets
% up the meanings of \cs{mfx@build@note}, \cs{mfx@build@skip}, and
% \cs{mfx@build@clear} for the pass; in the modes \texttt{off} and
% \texttt{warn} they are the ones defined below, and the mode
% \texttt{warn} first looks at the list without changing anything
% (section~\ref{sec:priority@down}).
%    \begin{macrocode}
  \csname MFX@down@\mfx@priority\endcsname
  \MFX@popdimen\Mfx@marginheight\mfx@pieceheights
%    \end{macrocode}
% Now we run forwards over the \cs{mfx@marginlist} to actually
% operate on each thing in the margin.
%    \begin{macrocode}
  \MFX@run@clear\mfx@marginlist
%<debug>\MFX@debug{buildmargin@down: RETURN
%<debug>           marginout=\MFX@mac\mfx@marginout,
%<debug>           marginlist=\MFX@mac\mfx@marginlist}%
}
%    \end{macrocode}
% \end{macros}
%
% We now define the various |\MFX@margin@...@down| macros.
% At this stage in the game, the only difference between  notes
% and skips is that we ignore skips before any notes by setting
% \cs{mfx@build@skip} initially to \cs{@gobble}.  Once we've
% seen the first note, skips are treated exactly the same: as
% fixed-height material.  If there is room in the current piece
% for the given height, then we prepend it to \cs{mfx@marginout},
% decrement the remaining height, and arrange for the boxes to
% be freed.  If not, we unshift the next piece height from
% \cs{mfx@pieceheights} and try again, until \cs{mfx@pieceheights}
% is empty and we simply defer everything to later pages.
%
% \begin{macros}{\MFX@margin@note@down}
% Upon seeing a note, we must do several things:
% \begin{enumerate}
% \item determine which box (left or right) is needed for the
%   current page, by calling \cs{MFX@whichbox}
% \item if the box fits, free both boxes, prepend \cs{mfx@marginout}
%   with a call to \cs{mfx@build@note}, and re-enable skips
% \item otherwise, defer the current note and all future notes
% \end{enumerate}
% The latter two steps are taken care of by \cs{MFX@margin@fit},
% which takes the height and two blocks of material: one to prepend
% to \cs{mfx@marginout} if it fits, the other to append to
% \cs{mfx@marginlist} if it doesn't.
%    \begin{macrocode}
\def\MFX@margin@note@down#1#2#3{%
%<debug>\MFX@debug{margin@note@down: ENTRY: #1/ #2 at #3}%
  \MFX@whichbox\@marbox#1#2%
  \if\MFX@check@fit{}{\ht\@marbox+\dp\@marbox}%
    \MFX@snoc\mfx@marginout{%
      \noexpand\@cons\noexpand\@freelist#1%
      \noexpand\@cons\noexpand\@freelist#2%
      \noexpand\mfx@build@note\@marbox{#3}}%
    \let\mfx@build@skip\MFX@margin@skip@down
  \else
    \mfx@build@clear
    \mfx@build@note{#1}{#2}{#3}%
  \fi
}
%    \end{macrocode}
% \end{macros}
%
% \begin{macros}{\MFX@margin@skip@down}
% Skips are similar.  A skip needs only to save itself back into
% \cs{mfx@marginout}, provided it fits.  If not, there is no need
% to defer it because it will just get gobbled at the top of the
% next page anyway.
%    \begin{macrocode}
\def\MFX@margin@skip@down#1{%
%<debug>\MFX@debug{margin@skip@down #1}%
  \if\MFX@check@fit{}{#1}%
    \MFX@snoc\mfx@marginout{\noexpand\mfx@build@skip{#1}}%
  \else
    \mfx@build@clear
  \fi
}
%    \end{macrocode}
% \end{macros}
%
% \begin{macros}{\MFX@margin@clear@down}
% Finally, \cs{MFX@margin@clear@down} is the only place we actually
% need to handle full-margin clears, since the downward pass does not
% ever push \cs{mfx@build@clear} onto \cs{mfx@marginout}.  When we
% see this, we simply redefine all three commands to append themselves
% back to \cs{mfx@marginlist}.
%    \begin{macrocode}
\def\MFX@build@clear@down{%
%<debug>\MFX@debug{clear@down}%
  \def\mfx@build@note##1##2##3{%
    \MFX@cons\mfx@marginlist{\noexpand\mfx@build@note##1##2{\MFX@minus@inf}}}%
  \def\mfx@build@skip##1{%
    \MFX@cons\mfx@marginlist{\noexpand\mfx@build@skip{##1}}}%
  \def\mfx@build@clear{%
    \MFX@cons\mfx@marginlist{\noexpand\mfx@build@clear}}%
}
%    \end{macrocode}
% \end{macros}
%
% \begin{macros}{\MFX@check@fit}
% We factored out some of the common functionality between the
% note and skip routines, so that must now be defined.  The
% \cs{MFX@check@fit} macro acts as a conditional and should be
% used as \cs{if}\cs{MFX@check@fit}\marg{piece-hook}\marg{size}.
% It takes care of iterating through the list of heights and
% accumulating the total size of material encountered so far.
% The \emph{piece-hook} is executed each time a new piece height
% is popped.
%    \begin{macrocode}
\def\MFX@check@fit#1#2{%
  00\fi % close out the \if
%<debug>\MFX@debug{check@fit{\unexpanded{#1}}{#2=\the\dimexpr#2} ENTRY:
%<debug>           marginheight=\the\Mfx@marginheight}%
  \@tempswafalse
  \ifdim\dimexpr#2<\Mfx@marginheight % it fits
    \advance\Mfx@marginheight-\dimexpr#2\relax % deduct the size
    \@tempswatrue
  \else % didn't fit: check the next piece
%<debug>\MFX@debug{check@fit overflow: pieceheights=\MFX@mac\mfx@pieceheights}%
    \ifx\mfx@pieceheights\@empty\else % make sure there's anything there
      #1%
      \MFX@popdimen\Mfx@marginheight\mfx@pieceheights
      \if\MFX@check@fit{#1}{#2}\fi
    \fi
  \fi
%<debug>\MFX@debug{check@fit RETURN \meaning\if@tempswa:
%<debug>           marginheight=\the\Mfx@marginheight,}%
  \if@tempswa % start a new \if
}
%    \end{macrocode}
% \end{macros}
% \begin{macros}{\MFX@popdimen}
% Here is a quick convenience routine.  \cs{MFX@popdimen}\marg{dimen}%
% \marg{list} removes the first dimension from \emph{list} and stores
% it into \emph{dimen}.
%    \begin{macrocode}
\def\MFX@popdimen#1#2{%
  \def\@elt##1{%
    #1##1\relax
    \def\@elt####1{%
      \MFX@cons#2{\noexpand\@elt{####1}}%
    }%
  }%
  \MFX@run@clear#2%
}
%    \end{macrocode}
% \end{macros}
%
% \begin{macros}{\MFX@whichbox}
% We also need to determine which box should be used, since they
% may have different heights.  The macro \cs{MFX@whichbox}%
% \marg{target-box}\marg{left-box}\marg{right-box} checks which margin
% we're setting and stores the correct box into \emph{target-box}.
% Note that \emph{target-box} must be a single control sequence.
%    \begin{macrocode}
\def\MFX@whichbox#1#2#3{%
  \if\MFX@leftmargin
    \def#1{#2}%
  \else
    \def#1{#3}%
  \fi
%<debug>\MFX@debug{whichbox: \@marbox (\the\dimexpr\ht#1+\dp#1)}%
}
%    \end{macrocode}
% \end{macros}
%
% \begin{macros}{\MFX@leftmargin}
% And here is the logic to figure out which margin we're in, based on
% the page number and other flags.  This is another conditional-like
% macro, and should be used after an \cs{if}, as in
% \cs{if}\cs{MFX@leftmargin}\ldots\cs{else}\ldots\cs{fi}.
%
% This is different from the corresponding code in the {\LaTeX}
% routines because we don't support double columns.  In addition, we
% would ideally allow \cs{if@reversemargin} to work on a per-note
% basis (i.e. at callout time) but we also need something working at
% shipout time so we can figure out which margin to use.  Thus, until
% we figure out how to use multiple margins, this will have to do.
%    \begin{macrocode}
\def\MFX@leftmargin{%
  00\fi % close out the \if
  \@tempcnta\@ne
  \if@mparswitch
    \unless\ifodd\c@page
      \@tempcnta\m@ne
    \fi
  \fi
  \if@reversemargin
    \@tempcnta-\@tempcnta
  \fi
%<debug>\MFX@debug{margin on \ifnum\@tempcnta<\z@ left\else right\fi}%
  \ifnum\@tempcnta<\z@ % start a new \if
}
%    \end{macrocode}
% \end{macros}
%
% \begin{macros}{\MFX@minus@inf}
% Finally, note that when deferring notes to the next page, we
% adjust their position to the top of the page, rather than the
% callout position.  This is a large negative dimension (near
% \TeX's maximum), but we may reconsider making this zero or even
% a small positive amount, since there seems to be a small amount
% of space before the first paragraph in normal text, though I'm
% not sure where that comes from.
%    \begin{macrocode}
\def\MFX@minus@inf{-4000\p@}
%    \end{macrocode}
% \end{macros}
%
% \begin{macros}{\MFX@down@off}
% The setup of the pass as described above, which is all the modes
% \texttt{off} and \texttt{warn} do.
%    \begin{macrocode}
\def\MFX@down@off{%
  \let\mfx@build@note\MFX@margin@note@down
  \let\mfx@build@skip\@gobble
  \let\mfx@build@clear\MFX@build@clear@down
  \let\mfx@marginout\@empty
}
%    \end{macrocode}
% \end{macros}
%
% \subsection{Priority in the first pass}
% \label{sec:priority@down}
% The pass above takes the notes strictly in order: it adds up their
% heights, and the first note that doesn't fit moves to the next page
% together with every note after it.  Where a note is called out plays no
% part in this; that only matters in the later passes, which place the notes
% that were taken.  So if deferrable notes early in the list use up the
% margin, an ordinary note later in the list moves, even though it would
% have fit had one of the deferrable notes moved instead.
%
% To find out whether that happens, we run the pass as a \emph{trial} first:
% the same fitting, in the same order, but without changing anything.  The
% trial reports the first ordinary note~$N$ that doesn't fit and the last
% deferrable note~$L$ that was taken before it.  In the mode \texttt{warn}
% the trial follows the order exactly like the real pass (a deferrable note
% that doesn't fit ends the page), and if it finds both $N$ and~$L$ we warn,
% naming $N$ and the deferrable notes ahead of it, and then run the real pass
% as always.  Hence that mode never changes the layout.
%
% In the modes \texttt{defer} and \texttt{split} we instead put $L$ on a
% \emph{skip list} and try again, until no ordinary note is kept off the page
% by a deferrable one.  In these trials a deferrable note that is skipped, or
% doesn't fit, takes all later deferrable notes with it to the next page, so
% that the deferrable notes keep their order, but the ordinary notes after
% it are still tried; the ordinary notes keep their order as before, since
% the first one that doesn't fit ends the page.  So the rule is: a
% deferrable note never holds up an ordinary note after it, but no note
% overtakes a note of its own kind, and nothing overtakes an ordinary note.
% This also fills the page after a long deferrable note that doesn't fit at
% all, which in the mode \texttt{warn} (as in the pass above) would take
% every later note to the next page.  We could let a later, shorter
% deferrable note fill the gap too, but commentary notes are often read in
% sequence, and a note carried to the next page ahead of the next page's
% own deferrable notes (\cs{MFX@carrylate}) relies on their order.
% A deferrable note that
% doesn't fit gives its space back: the fitting test moves on to the next
% piece of the margin when a note doesn't fit in the current one, and a note
% we skip must not use up the pieces.  Every trial skips one more note, so
% there are at most as many trials as there are deferrable notes on the page.
% Then the real pass takes the notes as the last trial did, moving the notes
% on the skip list, and the later deferrable ones, to the next page.
%
% In the mode \texttt{split}, the first deferrable note to move, whether
% skipped or not fitting, is split instead: the
% last trial also tells us how much room is left at the bottom of the
% margin, and we split the note to fit it, preferring a paragraph break
% (\cs{MFX@split} explains how).  The top part takes
% the note's place in the list, and the rest goes to the next page as a new
% deferrable note, with a new pair of registers from \cs{@freelist}.  We
% only split when the margin is a single piece, since then the room left
% over is room below the note; with several pieces it might be in another
% piece.
%
% A deferrable note moved whole goes to the top of the next page, like any
% note that didn't fit.  But there it would push that page's ordinary notes
% down, perhaps far from their marks, or even off that page too.  So in the
% modes \texttt{defer} and \texttt{split} we first move such notes, before
% the trials, to just before the first deferrable note of the page that
% wasn't moved (to keep the deferrable notes in order), or to the end.  A
% remainder of a split note stays at the top, where the reader continues.
%
% \begin{macros}{\Mfx@idx,\ifmfx@lowblocked,\ifmfx@full,\ifmfx@onepiece,
%     \Mfx@savedheight,\Mfx@splitroom,\Mfx@splitmin,\Mfx@split@tmp,
%     \Mfx@split@top}
% We number the notes in the list with \cs{Mfx@idx}, so that the trials and
% the real pass can refer to the same note.  \cs{ifmfx@lowblocked} is true
% once a deferrable note has moved and so all later ones must, and
% \cs{ifmfx@full} once the page is full.  \cs{Mfx@splitmin} is the least
% room in which we split a note.
%    \begin{macrocode}
\newcount\Mfx@idx
\newif\ifmfx@lowblocked
\newif\ifmfx@full
\newif\ifmfx@onepiece
\newdimen\Mfx@savedheight
\newdimen\Mfx@splitroom
\newdimen\Mfx@splitmin
\Mfx@splitmin=24\p@
\newbox\Mfx@split@tmp
\newbox\Mfx@split@top
%    \end{macrocode}
% \end{macros}
%
% \begin{macros}{\MFX@down@warn,\MFX@down@defer,\MFX@down@split}
% The three modes that look at the list.  The mode \texttt{split} differs
% from \texttt{defer} only in the real pass, where it tests the mode.
%    \begin{macrocode}
\def\MFX@down@warn{%
  \MFX@trials
  \MFX@down@off
}
\def\MFX@down@defer{%
  \MFX@carrylate
  \MFX@trials
  \Mfx@idx\z@
  \mfx@lowblockedfalse
  \mfx@fullfalse
  \let\mfx@build@note\MFX@real@note
  \let\mfx@build@skip\@gobble
  \let\mfx@build@clear\MFX@real@clear
  \let\mfx@marginout\@empty
}
\let\MFX@down@split\MFX@down@defer
\def\MFX@priority@warn{warn}
\def\MFX@priority@split{split}
%    \end{macrocode}
% \end{macros}
%
% \begin{macros}{\MFX@trials,\mfx@skipset}
% The trials start from the list of piece heights, which we keep in
% \cs{mfx@pieceheights@all} and put back after each trial, since the fitting
% test uses it up.  The skip list \cs{mfx@skipset} holds note numbers
% separated by commas.
%    \begin{macrocode}
\def\MFX@trials{%
  \let\mfx@pieceheights@all\mfx@pieceheights
  \begingroup
    \count@\z@
    \def\@elt##1{\advance\count@\@ne}%
    \mfx@pieceheights
    \ifnum\count@=\@ne
      \aftergroup\mfx@onepiecetrue
    \else
      \aftergroup\mfx@onepiecefalse
    \fi
  \endgroup
  \def\mfx@skipset{,}%
  \MFX@trials@loop
  \global\let\mfx@pieceheights\mfx@pieceheights@all
}
\def\MFX@trials@loop{%
  \MFX@trial
  \let\mfx@next\relax
  \ifnum\mfx@trial@L>\z@
    \ifx\mfx@priority\MFX@priority@warn
      \PackageWarningNoLine{marginfix}{Margin note from \mfx@trial@Nloc\space
        moved from page \thepage\MessageBreak to a later page; deferrable
        margin notes\MessageBreak ahead of it: \mfx@trial@lows
        \ifx\mfx@trial@blocker\@empty\else
          \MessageBreak (\mfx@trial@blocker\space did not fit and
          took\MessageBreak the notes after it along)%
        \fi
        \MessageBreak (shorten or move them, or use
        \string\marginnotepriority)}%
    \else
      \edef\mfx@skipset{\mfx@skipset\mfx@trial@L,}%
      \let\mfx@next\MFX@trials@loop
    \fi
  \fi
  \mfx@next
}
%    \end{macrocode}
% \end{macros}
%
% \begin{macros}{\MFX@trial,\MFX@trial@note,\MFX@trial@skip}
% One trial runs over the list in a group, so that everything but the
% results is forgotten afterwards; the results are \cs{mfx@trial@N} and
% \cs{mfx@trial@L} (zero if there is none), the location of~$N$, the
% locations of the deferrable notes taken before~$N$, the location of a
% deferrable note that didn't fit and ended the page (mode \texttt{warn}),
% and the room left.
%    \begin{macrocode}
\def\MFX@trial{%
  \begingroup
    \global\let\mfx@pieceheights\mfx@pieceheights@all
    \MFX@popdimen\Mfx@marginheight\mfx@pieceheights
    \Mfx@idx\z@
    \mfx@lowblockedfalse
    \mfx@fullfalse
    \gdef\mfx@trial@N{0}%
    \gdef\mfx@trial@L{0}%
    \def\mfx@lastlow{0}%
    \global\let\mfx@trial@lows\@empty
    \global\let\mfx@trial@blocker\@empty
    \let\mfx@build@note\MFX@trial@note
    \let\mfx@build@skip\@gobble
    \def\mfx@build@clear{\mfx@fulltrue}%
    \mfx@marginlist
    \global\Mfx@splitroom\Mfx@marginheight
  \endgroup
}
%    \end{macrocode}
% A deferrable note is taken if it isn't blocked and fits.  If it doesn't
% fit, the mode \texttt{warn} fills the page, like the real pass, and counts
% the note among those in the way; the other modes only block the later
% deferrable notes.
%    \begin{macrocode}
\def\MFX@trial@note#1#2#3{%
  \advance\Mfx@idx\@ne
  \let\mfx@build@skip\@gobble
  \ifnum\mfx@trial@N=\z@
    \MFX@whichbox\@marbox#1#2%
    \MFX@islow#1%
    \if@tempswa
      \ifmfx@full\else
        \MFX@inskip
        \ifin@\mfx@lowblockedtrue\fi
        \ifmfx@lowblocked\else
          \MFX@try@fit{\ht\@marbox+\dp\@marbox}%
          \if@tempswa
            \edef\mfx@lastlow{\the\Mfx@idx}%
            \xdef\mfx@trial@lows{\mfx@trial@lows\MFX@loc#1\space}%
            \let\mfx@build@skip\MFX@trial@skip
          \else
            \ifx\mfx@priority\MFX@priority@warn
              \mfx@fulltrue
              \edef\mfx@lastlow{\the\Mfx@idx}%
              \xdef\mfx@trial@lows{\mfx@trial@lows\MFX@loc#1\space}%
              \xdef\mfx@trial@blocker{\MFX@loc#1}%
            \else
              \mfx@lowblockedtrue
            \fi
          \fi
        \fi
      \fi
    \else
%    \end{macrocode}
% An ordinary note that doesn't fit, or comes after the page is full, is our
% $N$.
%    \begin{macrocode}
      \@tempswafalse
      \ifmfx@full\else
        \if\MFX@check@fit{}{\ht\@marbox+\dp\@marbox}%
          \let\mfx@build@skip\MFX@trial@skip
        \fi
      \fi
      \if@tempswa\else
        \xdef\mfx@trial@N{\the\Mfx@idx}%
        \xdef\mfx@trial@L{\mfx@lastlow}%
        \xdef\mfx@trial@Nloc{\MFX@loc#1}%
      \fi
    \fi
  \fi
}
\def\MFX@trial@skip#1{%
  \if\MFX@check@fit{}{#1}\else
    \mfx@fulltrue
  \fi
}
%    \end{macrocode}
% \end{macros}
%
% \begin{macros}{\MFX@try@fit,\MFX@inskip}
% \cs{MFX@try@fit}\marg{size} is \cs{MFX@check@fit} for a deferrable note:
% if the note doesn't fit, it puts back the space and pieces it used up.  It
% leaves the result in \cs{@tempswa}.  \cs{MFX@inskip} sets \cs{ifin@} to
% whether the current note is on the skip list.
%    \begin{macrocode}
\def\MFX@try@fit#1{%
  \Mfx@savedheight\Mfx@marginheight
  \let\mfx@saved@pieceheights\mfx@pieceheights
  \if\MFX@check@fit{}{#1}%
  \else
    \Mfx@marginheight\Mfx@savedheight
    \global\let\mfx@pieceheights\mfx@saved@pieceheights
  \fi
}
\def\MFX@inskip{%
  \edef\reserved@a{\noexpand\in@{,\the\Mfx@idx,}{\mfx@skipset}}%
  \reserved@a
}
%    \end{macrocode}
% \end{macros}
%
% \begin{macros}{\MFX@real@note}
% The real pass of the modes \texttt{defer} and \texttt{split}.  It does
% what \cs{MFX@margin@note@down} does, but takes and moves notes as the last
% trial did.  A note on the skip list is split, if we may, or else moved;
% either way we warn.  So is a deferrable note that doesn't fit, but then we
% only warn if we split it, since marginfix has always moved such notes.
% Only the first deferrable note to move gets this far: the deferrable
% notes after it are blocked already.
%    \begin{macrocode}
\def\MFX@real@note#1#2#3{%
  \advance\Mfx@idx\@ne
  \MFX@whichbox\@marbox#1#2%
  \ifmfx@full
    \MFX@defer{#1}{#2}%
  \else
    \MFX@islow#1%
    \if@tempswa
      \ifmfx@lowblocked\in@false\else\MFX@inskip\fi
      \ifin@
        \MFX@give@way{#1}{#2}{#3}\@tempswatrue
      \else
        \ifmfx@lowblocked
          \MFX@defer{#1}{#2}%
        \else
          \MFX@try@fit{\ht\@marbox+\dp\@marbox}%
          \if@tempswa
            \MFX@place{#1}{#2}{#3}%
          \else
            \MFX@give@way{#1}{#2}{#3}\@tempswafalse
          \fi
        \fi
      \fi
    \else
      \if\MFX@check@fit{}{\ht\@marbox+\dp\@marbox}%
        \MFX@place{#1}{#2}{#3}%
      \else
        \mfx@fulltrue
        \MFX@defer{#1}{#2}%
      \fi
    \fi
  \fi
}
%    \end{macrocode}
% \end{macros}
%
% \begin{macros}{\MFX@give@way,\ifmfx@warnmove}
% \cs{MFX@give@way}\marg{box}\marg{box}\marg{position}\marg{warn} moves the
% first deferrable note that must move, splitting it in the mode
% \texttt{split} if it can.  The last argument sets \cs{@tempswa} to whether
% to warn when the note moves whole.  The reason given in the warnings is
% the same either way, since a note that doesn't fit leaves its room to the
% notes after it too.
%    \begin{macrocode}
\newif\ifmfx@warnmove
\def\MFX@give@way#1#2#3#4{%
  \mfx@lowblockedtrue
  \edef\mfx@thisloc{\MFX@loc#1}%
  #4%
  \if@tempswa\mfx@warnmovetrue\else\mfx@warnmovefalse\fi
  \@tempswafalse
  \ifx\mfx@priority\MFX@priority@split
    \ifmfx@onepiece
      \MFX@split{#1}{#2}{#3}%
    \fi
  \fi
  \if@tempswa
    \PackageWarningNoLine{marginfix}{Deferrable margin note from
      \mfx@thisloc\MessageBreak split at the end of page \thepage,
      to keep the\MessageBreak margin notes after it on their page}%
  \else
    \ifmfx@warnmove
      \PackageWarningNoLine{marginfix}{Deferrable margin note from
        \mfx@thisloc\MessageBreak moved from page \thepage\space
        to the next, to keep\MessageBreak the margin notes after it
        on their page}%
    \fi
    \MFX@defer{#1}{#2}%
  \fi
}
%    \end{macrocode}
% \end{macros}
%
% \begin{macros}{\MFX@place,\MFX@defer,\MFX@real@skip,\MFX@real@clear}
% Taking a note is as in \cs{MFX@margin@note@down}, and moving one as in
% \cs{MFX@build@clear@down}, but a deferrable note moved whole is flagged
% for \cs{MFX@carrylate}.  A skip is taken if it follows a note that was
% taken and fits, and moved if it follows a note that was moved.  The first
% \cs{clearmargin} fills the page, and later ones are kept for the next.
%    \begin{macrocode}
\def\MFX@place#1#2#3{%
  \expandafter\global\expandafter\let
    \csname mfx@carried@\number#1\endcsname\relax
  \MFX@snoc\mfx@marginout{%
    \noexpand\@cons\noexpand\@freelist#1%
    \noexpand\@cons\noexpand\@freelist#2%
    \noexpand\mfx@build@note\@marbox{#3}}%
  \let\mfx@build@skip\MFX@real@skip
}
\def\MFX@defer#1#2{%
  \MFX@islow#1%
  \if@tempswa
    \expandafter\gdef\csname mfx@carried@\number#1\endcsname{}%
  \fi
  \MFX@cons\mfx@marginlist{\noexpand\mfx@build@note#1#2{\MFX@minus@inf}}%
  \def\mfx@build@skip##1{%
    \MFX@cons\mfx@marginlist{\noexpand\mfx@build@skip{##1}}}%
}
\def\MFX@real@skip#1{%
  \if\MFX@check@fit{}{#1}%
    \MFX@snoc\mfx@marginout{\noexpand\mfx@build@skip{#1}}%
  \else
    \mfx@fulltrue
    \let\mfx@build@skip\@gobble
  \fi
}
\def\MFX@real@clear{%
  \ifmfx@full
    \MFX@cons\mfx@marginlist{\noexpand\mfx@build@clear}%
  \else
    \mfx@fulltrue
  \fi
  \let\mfx@build@skip\@gobble
}
%    \end{macrocode}
% \end{macros}
%
% \begin{macros}{\MFX@split}
% \cs{MFX@split}\marg{box}\marg{box}\marg{position} splits a note to the room
% left by the last trial, less the skip after the note, and sets
% \cs{@tempswa} if it did.
%
% We don't leave the choice of the break to \cs{vsplit}.  \cs{vsplit} weighs
% each break by the badness of the top part plus the penalty there, but a
% margin note has no stretchable glue, so every top part that doesn't fill
% the room exactly is infinitely bad, all breaks cost the same, and
% \cs{vsplit} takes the last one that fits, whatever its penalty.  So it
% splits in the middle of a sentence, or of a citation, as happily as
% between paragraphs, and ignores the club, widow and interline penalties
% of the note's paragraphs.
%
% Instead we look at the list ourselves.  A note's box holds a \cs{vtop}
% (see \cs{@savemarbox}); we take it out with \cs{lastbox} and then take a
% copy of its list apart from the bottom, with \cs{lastbox}, \cs{unskip},
% \cs{unkern} and \cs{unpenalty}, recognising each item by e-\TeX's
% \cs{lastnodetype} and keeping track of the height of what is left above
% it.  Between two lines there is a \emph{gap} of glue and penalties, and
% the paragraph builder put the club, widow, interline and broken penalties
% of the note's paragraphs there; a gap is a legal break if it starts with
% glue (at no cost) or has a penalty below 10000 (at that cost).  Between
% two paragraphs the gap holds the \cs{parskip} glue as well as the
% interline glue, so a gap with at least two pieces of glue is a paragraph
% break.  (\cs{parskip} is usually zero in a note, so we can't tell it by
% its size.)  We can't take apart anything else, such as a \cs{write} from
% a \cs{label} between paragraphs, so the scan stops there, and breaks
% above it are not considered.
%
% Of the breaks where the top part fits, we choose
% \begin{enumerate}
% \item the last paragraph break that leaves at most the \emph{slack}
%   unused, where the slack is \cs{marginsplitslack} but at most half the
%   room;
% \item otherwise the cheapest line break within the slack, the last of
%   them if several cost the same, so that with the usual penalties we
%   don't leave a single line of a paragraph on either side;
% \item otherwise the last break that fits at all;
% \item and if the scan found none, because it stopped early, the break
%   \cs{vsplit} would take.
% \end{enumerate}
% The slack bounds what we give up for a paragraph break: the commentary
% that moves to the next page grows by at most that much, and it pushes the
% next page's notes down by as much.  A few lines is a fair price for not
% cutting a thought in half, and the bound by half the room keeps us from
% leaving a nearly empty margin when little room is left.  We don't try to
% avoid breaking inside a sentence: \TeX{} can't see the characters at the
% end of a line (Lua\TeX{} could), and a heuristic on the glue wouldn't be
% reliable.  The paragraph breaks, and the penalties an author can add, are
% what we have.
%
% Having chosen, we let \cs{vsplit} make the cut, to the height of the top
% part we measured: without stretch it takes the last break that fits, and
% that is ours, since the next break down is higher by at least a line.
% \cs{vsplit} leaves whatever is in the list, \cs{write}s included, in the
% part where it was, and drops the glue and penalties at the break, as at a
% page break.
%
% The rest starts with the cue \cs{marginnotecontinued}, on a line of its
% own, in the font the note had where it said \cs{marginnotedeferrable}; we
% can't run the cue into the first line, since that is typeset already.  We get
% the interline glue after the cue from \cs{vsplit} by setting
% \cs{splittopskip} to the distance from the cue's depth to the next
% baseline.  Without a cue, \cs{splittopskip} is the height of the rest's
% first line, which the scan noted, so that the rest starts with a line, as
% the note did.  Either way the rest is a \cs{vtop} like the note, and is
% placed on the next page like any note moved to the top of the margin,
% with its first baseline at \cs{topskip}.  The top part can end with the
% cue \cs{marginnotecontinues}, set flush right, for which we reserve a
% line.
%
% \begin{macros}{\Mfx@split@rest,\Mfx@scan@size,\Mfx@scan@room,
%     \Mfx@scan@slack,\Mfx@scan@ht,\Mfx@gap@ht,\Mfx@gap@glue,\Mfx@gap@pen,
%     \Mfx@best@pen,\Mfx@cut@size,\Mfx@cut@ht,\Mfx@par@size,\Mfx@par@ht,
%     \Mfx@line@size,\Mfx@line@ht,\Mfx@any@size,\Mfx@any@ht,
%     \ifmfx@scan@box,\ifmfx@scan@ingap}
% The state of the scan: the height of what is left (\cs{Mfx@scan@size}),
% the height of the last box taken off (\cs{Mfx@scan@ht}), and for the
% current gap the height of the box below it, how much glue it has and its
% lowest penalty.  A break is recorded by the height of the top part and
% the height of the first line of the rest.
%    \begin{macrocode}
\newbox\Mfx@split@rest
\newdimen\Mfx@scan@size
\newdimen\Mfx@scan@room
\newdimen\Mfx@scan@slack
\newdimen\Mfx@scan@ht
\newdimen\Mfx@gap@ht
\newcount\Mfx@gap@glue
\newcount\Mfx@gap@pen
\newcount\Mfx@best@pen
\newdimen\Mfx@cut@size
\newdimen\Mfx@cut@ht
\newdimen\Mfx@par@size
\newdimen\Mfx@par@ht
\newdimen\Mfx@line@size
\newdimen\Mfx@line@ht
\newdimen\Mfx@any@size
\newdimen\Mfx@any@ht
\newif\ifmfx@scan@box
\newif\ifmfx@scan@ingap
%    \end{macrocode}
% \end{macros}
%
% First the split itself.  We work in a group, in the note's font, and take
% the list out of the note's box.
%    \begin{macrocode}
\def\MFX@split#1#2#3{%
  \begingroup
    \global\let\mfx@split@ok\@empty
    \MFX@notefont#1%
    \Mfx@scan@room\dimexpr\Mfx@splitroom-\marginparpush-\p@\relax
    \ifx\mfx@contsuffix\@empty\else
      \advance\Mfx@scan@room-\baselineskip
    \fi
    \ifdim\Mfx@scan@room>\Mfx@splitmin
      \setbox\z@\vbox{\unvcopy\@marbox\global\setbox\Mfx@split@tmp\lastbox}%
      \ifvoid\Mfx@split@tmp\else
        \setbox\Mfx@split@tmp\vbox{\unvbox\Mfx@split@tmp}%
        \MFX@split@choose
        \MFX@split@cut
      \fi
    \fi
  \endgroup
%    \end{macrocode}
% If the split worked we need two free registers for the rest.  If we get
% them, the top part replaces the note in both of its registers and is taken
% (it fits, but we test it the usual way to use up the room), and the rest
% goes to the next page right after it, as a deferrable note with the same
% font for its cues, should it be split again.
%    \begin{macrocode}
  \@tempswafalse
  \ifx\mfx@split@ok\@empty\else
    \@next\mfx@restA\@freelist{%
      \@next\mfx@restB\@freelist{\@tempswatrue}%
        {\@cons\@freelist\mfx@restA}}{}%
  \fi
  \if@tempswa
    \global\setbox#1\copy\Mfx@split@top
    \global\setbox#2\box\Mfx@split@top
    \global\setbox\mfx@restA\copy\Mfx@split@rest
    \global\setbox\mfx@restB\box\Mfx@split@rest
    \MFX@forget\mfx@restA
    \MFX@forget\mfx@restB
    \MFX@copyclass\mfx@restA#1%
    \MFX@copyclass\mfx@restB#1%
    \expandafter\xdef\csname mfx@boxloc@\number\mfx@restA\endcsname{%
      \MFX@loc#1 (continued)}%
    \MFX@try@fit{\ht#1+\dp#1}%
    \if@tempswa
      \MFX@place{#1}{#2}{#3}%
    \else
      \MFX@defer{#1}{#2}%
    \fi
    \MFX@cons\mfx@marginlist{%
      \noexpand\mfx@build@note\mfx@restA\mfx@restB{\MFX@minus@inf}%
      \noexpand\mfx@build@skip{\the\marginparpush}}%
    \@tempswatrue
  \fi
}
%    \end{macrocode}
% \cs{MFX@copyclass}\marg{to}\marg{from} marks the note in box \emph{to}
% deferrable, with the font of the note in box \emph{from}.
%    \begin{macrocode}
\def\MFX@copyclass#1#2{%
  \expandafter\gdef\csname mfx@cls@\number#1\endcsname{}%
  \expandafter\global\expandafter\let
    \csname mfx@nf@\number#1\expandafter\endcsname
    \csname mfx@nf@\number#2\endcsname
  \expandafter\global\expandafter\let
    \csname mfx@nb@\number#1\expandafter\endcsname
    \csname mfx@nb@\number#2\endcsname
}
%    \end{macrocode}
% Choosing the break: the slack, then the scan, in a box we throw away, and
% then the choice in order of preference, leaving the height of the top
% part and of the rest's first line in \cs{Mfx@cut@size} and
% \cs{Mfx@cut@ht}.  If the scan found nothing, we cut at the room and guess
% the first line to be as high as a strut.
%    \begin{macrocode}
\def\MFX@split@choose{%
  \Mfx@scan@slack\marginsplitslack
  \ifdim\Mfx@scan@slack>.5\Mfx@scan@room
    \Mfx@scan@slack.5\Mfx@scan@room
  \fi
  \global\Mfx@par@size\z@
  \global\Mfx@line@size\z@
  \global\Mfx@any@size\z@
  \global\Mfx@best@pen\@M
  \Mfx@scan@size\dimexpr\ht\Mfx@split@tmp+\dp\Mfx@split@tmp\relax
  \setbox\z@\vbox{\unvcopy\Mfx@split@tmp\MFX@scan}%
  \ifdim\Mfx@par@size>\z@
    \Mfx@cut@size\Mfx@par@size
    \Mfx@cut@ht\Mfx@par@ht
  \else\ifdim\Mfx@line@size>\z@
    \Mfx@cut@size\Mfx@line@size
    \Mfx@cut@ht\Mfx@line@ht
  \else\ifdim\Mfx@any@size>\z@
    \Mfx@cut@size\Mfx@any@size
    \Mfx@cut@ht\Mfx@any@ht
  \else
    \Mfx@cut@size\Mfx@scan@room
    \Mfx@cut@ht\ht\strutbox
  \fi\fi\fi
}
%    \end{macrocode}
% Cutting, and putting both parts into \cs{vtop}s, with their cues.  The
% top part must not be empty, must fit, and must leave a rest.  Without a
% cue, the rest begins with the (zero) glue from \cs{splittopskip}, so
% \cs{vtop} would give it no height; we set its height to that of its first
% line, which puts the first baseline at the reference point, as in the
% note.
%    \begin{macrocode}
\def\MFX@split@cut{%
  \splitmaxdepth\maxdimen
  \vbadness\@M
  \vfuzz\maxdimen
  \ifx\mfx@contprefix\@empty
    \splittopskip\Mfx@cut@ht
  \else
    \setbox\tw@\hbox{\mfx@contprefix}%
    \splittopskip\dimexpr\baselineskip-\dp\tw@\relax
  \fi
  \setbox\Mfx@split@top\vsplit\Mfx@split@tmp to \Mfx@cut@size
  \setbox\Mfx@split@top\vbox{\unvbox\Mfx@split@top}%
  \dimen@\dimexpr\ht\Mfx@split@top+\dp\Mfx@split@top\relax
  \ifdim\dimen@>\z@
    \ifdim\dimen@>\Mfx@scan@room\else
      \ifvoid\Mfx@split@tmp\else
        \dimen@\dp\Mfx@split@top
        \dimen@ii\wd\Mfx@split@top
        \global\setbox\Mfx@split@top\vbox{\vtop{%
          \unvbox\Mfx@split@top
          \ifx\mfx@contsuffix\@empty\else
            \prevdepth\dimen@
            \hbox to\dimen@ii{\hfil\mfx@contsuffix}%
          \fi}}%
        \ifx\mfx@contprefix\@empty
          \setbox\z@\vtop{\unvbox\Mfx@split@tmp}%
          \dp\z@\dimexpr\dp\z@-\Mfx@cut@ht\relax
          \ht\z@\Mfx@cut@ht
        \else
          \setbox\z@\vtop{\box\tw@\unvbox\Mfx@split@tmp}%
        \fi
        \global\setbox\Mfx@split@rest\vbox{\box\z@}%
        \gdef\mfx@split@ok{ok}%
      \fi
    \fi
  \fi
}
%    \end{macrocode}
% \end{macros}
%
% \begin{macros}{\MFX@scan,\MFX@scan@step}
% The scan runs inside the box being built, where \cs{lastbox} and friends
% work.  It takes items off until the list is empty, or it meets an item it
% can't take off, or \cs{MFX@scan@endgap} finds that the breaks from there
% on leave more than the slack unused.
%    \begin{macrocode}
\def\MFX@scan{%
  \mfx@scan@boxfalse
  \mfx@scan@ingapfalse
  \MFX@scan@step
}
\def\MFX@scan@step{%
  \let\mfx@scan@next\MFX@scan@step
  \ifcase\lastnodetype
    \let\mfx@scan@next\relax % 0: a character
  \or\MFX@scan@box           % 1: an hbox
  \or\MFX@scan@box           % 2: a vbox
  \or\let\mfx@scan@next\relax\or\let\mfx@scan@next\relax
  \or\let\mfx@scan@next\relax\or\let\mfx@scan@next\relax
  \or\let\mfx@scan@next\relax\or\let\mfx@scan@next\relax
  \or\let\mfx@scan@next\relax\or\let\mfx@scan@next\relax
  \or\MFX@scan@glue          % 11
  \or\MFX@scan@kern          % 12
  \or\MFX@scan@penalty       % 13
  \else\let\mfx@scan@next\relax % -1: empty, or anything else
  \fi
  \mfx@scan@next
}
%    \end{macrocode}
% \end{macros}
%
% \begin{macros}{\MFX@scan@box,\MFX@scan@glue,\MFX@scan@kern,
%     \MFX@scan@penalty,\MFX@scan@gapitem}
% Taking off one item.  A box ends the gap below it, if any.  The other
% items belong to a gap if there is a box below them.
%    \begin{macrocode}
\def\MFX@scan@box{%
  \ifmfx@scan@ingap\MFX@scan@endgap\fi
  \setbox\z@\lastbox
  \advance\Mfx@scan@size-\dimexpr\ht\z@+\dp\z@\relax
  \Mfx@scan@ht\ht\z@
  \mfx@scan@boxtrue
}
\def\MFX@scan@glue{%
  \MFX@scan@gapitem
  \advance\Mfx@scan@size-\lastskip
  \unskip
  \advance\Mfx@gap@glue\@ne
  \let\mfx@gap@top\MFX@scan@glue
}
\def\MFX@scan@kern{%
  \MFX@scan@gapitem
  \advance\Mfx@scan@size-\lastkern
  \unkern
  \let\mfx@gap@top\MFX@scan@kern
}
\def\MFX@scan@penalty{%
  \MFX@scan@gapitem
  \ifnum\lastpenalty<\Mfx@gap@pen \Mfx@gap@pen\lastpenalty \fi
  \unpenalty
  \let\mfx@gap@top\MFX@scan@penalty
}
\def\MFX@scan@gapitem{%
  \ifmfx@scan@box
    \ifmfx@scan@ingap\else
      \mfx@scan@ingaptrue
      \Mfx@gap@ht\Mfx@scan@ht
      \Mfx@gap@glue\z@
      \Mfx@gap@pen\@M
    \fi
  \fi
}
%    \end{macrocode}
% \end{macros}
%
% \begin{macros}{\MFX@scan@endgap}
% At the top of a gap we weigh the break there, as described above.  The
% size left is the height of the top part.  Scanning upwards, the breaks
% come in the order last to first, so the first paragraph break within the
% slack is the last one, and so on.  Once a break leaves more than the
% slack unused and we have one that fits, no break further up can be
% better, and we stop.
%    \begin{macrocode}
\def\MFX@scan@endgap{%
  \mfx@scan@ingapfalse
  \ifx\mfx@gap@top\MFX@scan@glue
    \count@\z@
  \else
    \count@\Mfx@gap@pen
  \fi
  \ifnum\count@<\@M
    \ifdim\Mfx@scan@size>\z@
      \ifdim\Mfx@scan@size>\Mfx@scan@room\else
        \ifdim\Mfx@any@size=\z@
          \global\Mfx@any@size\Mfx@scan@size
          \global\Mfx@any@ht\Mfx@gap@ht
        \fi
        \ifdim\dimexpr\Mfx@scan@room-\Mfx@scan@size\relax>\Mfx@scan@slack
          \let\mfx@scan@next\relax
        \else
          \ifnum\Mfx@gap@glue>\@ne
            \ifdim\Mfx@par@size=\z@
              \global\Mfx@par@size\Mfx@scan@size
              \global\Mfx@par@ht\Mfx@gap@ht
            \fi
          \fi
          \ifnum\count@<\Mfx@best@pen
            \global\Mfx@best@pen\count@
            \global\Mfx@line@size\Mfx@scan@size
            \global\Mfx@line@ht\Mfx@gap@ht
          \fi
        \fi
      \fi
    \fi
  \fi
}
%    \end{macrocode}
% \end{macros}
%
% \begin{macros}{\MFX@carrylate,\MFX@cl@flush}
% Before the trials, \cs{MFX@carrylate} moves the deferrable notes that were
% moved whole from the previous page to just before the first deferrable
% note that wasn't, or to the end of the list, each with the skip after it.
% We collect them in \cs{mfx@cl@carried} and flush them into the new list
% \cs{mfx@cl@out} when we meet a deferrable note or a \cs{clearmargin}.
%    \begin{macrocode}
\def\MFX@carrylate{%
  \begingroup
    \global\let\mfx@cl@carried\@empty
    \global\let\mfx@cl@out\@empty
    \def\mfx@cl@add{\MFX@cons\mfx@cl@out}%
    \def\mfx@build@note##1##2##3{%
      \expandafter\ifx\csname mfx@carried@\number##1\endcsname\relax
        \MFX@islow##1%
        \if@tempswa\MFX@cl@flush\fi
        \def\mfx@cl@add{\MFX@cons\mfx@cl@out}%
      \else
        \def\mfx@cl@add{\MFX@cons\mfx@cl@carried}%
      \fi
      \mfx@cl@add{\noexpand\mfx@build@note##1##2{##3}}}%
    \def\mfx@build@skip##1{\mfx@cl@add{\noexpand\mfx@build@skip{##1}}}%
    \def\mfx@build@clear{%
      \MFX@cl@flush
      \def\mfx@cl@add{\MFX@cons\mfx@cl@out}%
      \mfx@cl@add{\noexpand\mfx@build@clear}}%
    \mfx@marginlist
    \MFX@cl@flush
  \endgroup
  \global\let\mfx@marginlist\mfx@cl@out
}
\def\MFX@cl@flush{%
  \xdef\mfx@cl@out{\unexpanded\expandafter{\mfx@cl@out}%
    \unexpanded\expandafter{\mfx@cl@carried}}%
  \global\let\mfx@cl@carried\@empty
}
%    \end{macrocode}
% \end{macros}
%
% \subsection{Second pass: global upward}
% \begin{macros}{\MFX@buildmargin@up,\mfx@phantomheights}
% The next step is the global ``up'' step, in which we figure out the
% lowest piece a note can possibly occupy (without pushing later notes
% off the bottom) and add this information to the \cs{mfx@build@note} in
% \cs{mfx@marginout}.  We start similar to \cs{MFX@buildmargin@down},
% except we need the list of heights to be backwards.  We also need a
% list of phantom heights, in order to handle skips properly, which we
% intersperse as negative heights.
%    \begin{macrocode}
\def\MFX@buildmargin@up{%
%<debug>\MFX@debug{buildmargin@up: ENTRY
%<debug>           marginpieces=\MFX@mac\mfx@marginpieces,
%<debug>           marginout=\MFX@mac\mfx@marginout}%
  \let\mfx@pieceheights\@empty
  \let\mfx@phantomheights\@empty
  \let\temp@@\relax
  \def\@elt##1##2{%
%<debug>\MFX@debug{  -> piece (##1,##2), temp@@=\meaning\temp@@}%
    \MFX@snoc\mfx@pieceheights{\noexpand\@elt{\the\dimexpr##2-##1}}%
    \ifx\temp@@\relax\else
%<debug>\MFX@debug{  -> phantom (\temp@@,##1)}%
      \MFX@snoc\mfx@phantomheights{\noexpand\@elt{\the\dimexpr##1-\temp@@}}%
    \fi
    \def\temp@@{##2}%
  }%
  \mfx@in@phantomfalse
  \mfx@marginpieces
%    \end{macrocode}
% The piece counter, \cs{Mfx@piece@count} has not been touched yet,
% so now we will start decrementing it for each piece to keep track
% of where we are.  Note that \cs{mfx@marginout} will never contain
% \cs{mfx@build@clear} so we don't need to reassign it.  Since we
% don't do any deferrals here, we don't need to empty out a new target
% list.  Instead, we operate ``in-place'' in \cs{mfx@marginout}, after
% popping off the first height.
%    \begin{macrocode}
  \MFX@popdimen\Mfx@marginheight\mfx@pieceheights
  \let\mfx@build@note\MFX@margin@note@up
  \let\mfx@build@skip\@gobble
  \MFX@run@clear\mfx@marginout
%<debug>\MFX@debug{buildmargin@up: RETURN marginout=\MFX@mac\mfx@marginout}%
}
%    \end{macrocode}
% \end{macros}
%
% \begin{macros}{\MFX@margin@note@up}
% We must again define the specific behavior of each build command.
% These macros simply reuse \cs{MFX@check@fit}, but ask it to
% decrement the piece counter when a piece runs out of space.
% Aside from that, the only thing that actually happens here is
% that we append the current piece to each note, and also end
% up reversing the contents by again prepending everything.
% We are guaranteed that \cs{MFX@check@fit} will never fail.
% Since we cannot put notes in a phantom, we start by ensuring
% we're not in one.
%    \begin{macrocode}
\def\MFX@margin@note@up#1#2{%
%<debug>\MFX@debug{margin@note@up: #1at #2, marginheight=\the\Mfx@marginheight}%
  \ifmfx@in@phantom
    \MFX@popdimen\Mfx@marginheight\mfx@pieceheights
    \advance\Mfx@piece@count\m@ne
    \mfx@in@phantomfalse
  \fi
%    \end{macrocode}
% Now we're guaranteed to be in a piece, rather than a phantom, we
% look for a piece that can fit this note, making sure to decrement
% the piece count and pop off a phantom for each new piece we check.
% Once it's found, we add the note back to \cs{mfx@marginout} with
% the correct piece.
%    \begin{macrocode}
  \if\MFX@check@fit{\advance\Mfx@piece@count\m@ne
                    \MFX@popdimen\dimen@\mfx@phantomheights}{\ht#1+\dp#1}%
    \MFX@snoc\mfx@marginout{%
      \noexpand\mfx@build@note{#1}{#2}{\the\Mfx@piece@count}}%
    \let\mfx@build@skip\MFX@margin@skip@up
  \else\MFX@AssertionError\fi
}
%    \end{macrocode}
% \end{macros}
%
% \begin{macros}{\MFX@margin@skip@up}
% Skips are similar, but we have the added complication of handling
% margin phantoms.  When we cross between phantom and piece, we split
% the skip so that we can use the simplest recursion possible.
%    \begin{macrocode}
\def\MFX@margin@skip@up#1{%
%<debug>\MFX@debug{margin@skip@up: #1}%
  \dimen@#1\relax
  \advance\Mfx@marginheight-\dimen@
  \ifdim\Mfx@marginheight<\z@
%    \end{macrocode}
% This skip was bigger than the piece, so we need to split this skip,
% adding the overflow back, and try again.  Since we're done with this
% piece, we'll pop the next one and recurse on whatever's left.  This
% looks slightly different depending on whether or not we're in a phantom.
%    \begin{macrocode}
    \advance\dimen@\Mfx@marginheight
    \MFX@snoc\mfx@marginout{%
      \noexpand\mfx@build@skip{\the\dimen@}{\the\Mfx@piece@count}}%
    \dimen@-\Mfx@marginheight
    \ifmfx@in@phantom
      \MFX@popdimen\Mfx@marginheight\mfx@pieceheights
      \advance\Mfx@piece@count\m@ne
      \mfx@in@phantomfalse
    \else
      \MFX@popdimen\Mfx@marginheight\mfx@phantomheights
      \mfx@in@phantomtrue
    \fi
    \mfx@build@skip\dimen@
  \else
%    \end{macrocode}
% This skip fit entirely within the phantom, so we simply emit it.
%    \begin{macrocode}
    \MFX@snoc\mfx@marginout{%
      \noexpand\mfx@build@skip{\the\dimen@}{\the\Mfx@piece@count}}%
  \fi
}
%    \end{macrocode}
% \end{macros}
%
% \subsection{Margin pieces}
% \begin{macros}{\MFX@buildmargin@pieces}
% Before we can start the third and fourth passes, we need to set up
% a loop over the pieces so that each piece can do these passes at
% one time.  In case we didn't use up all the pieces in the second
% phase, we'll reset \cs{Mfx@piece@count} to zero.  We also reset
% \cs{Mfx@marginspace} and \cs{Mfx@marginbox}.
%    \begin{macrocode}
\def\MFX@buildmargin@pieces{%
  \Mfx@piece@count\z@
  \Mfx@marginspace\z@
  \setbox\Mfx@marginbox\vbox{\vskip\z@}%  TODO - do we need this?
%    \end{macrocode}
% Now we run over the individual margin pieces, clearing it out as we
% build up the contents of \cs{Mfx@marginbox}.  Once we're done with
% that, we need to do a bit of clean-up before finishing.
%    \begin{macrocode}
  \let\@elt\MFX@buildmargin@piece
  \MFX@run@clear\mfx@marginpieces
  \let\@elt\relax
  \global\Mfx@piece@count\z@
}
%    \end{macrocode}
% \end{macros}
%
% \begin{macros}{\MFX@buildmargin@piece}
% The \cs{MFX@buildmargin@piece} macro is called for each piece of
% the margin, and is passed the top and bottom positions of the piece.
% Here we need to do a few things.  First, if the output so far is
% smaller than the top of the piece (this is generally true, except
% sometimes for the first piece) then we need to insert padding into
% \cs{Mfx@marginbox} before continuing.  We accumulate this padding
% into \cs{Mfx@marginspace}, which allows skips to do double-duty
% across phantoms.
%    \begin{macrocode}
\def\MFX@buildmargin@piece#1#2{%
%<debug>\MFX@debug{buildmargin@piece: ENTRY (#1, #2)}%
  \ifdim\ht\Mfx@marginbox<#1\relax
    \dimen@\dimexpr#1-\ht\Mfx@marginbox\relax
%<debug>\MFX@debug{buildmargin@piece: padding \the\dimen@}%
    \setbox\Mfx@marginbox\vbox{%
      \unvbox\Mfx@marginbox
      \vskip\dimen@
    }%
    \advance\Mfx@marginspace\dimen@
  \fi
%    \end{macrocode}
% Now that \cs{Mfx@marginbox} has been padded, we proceed to set
% things up for our down and up passes, and then run them to build
% \cs{Mfx@piece@content}.
%    \begin{macrocode}
  \Mfx@marginpos@min#1\relax
  \Mfx@marginpos@max#1\relax
  \Mfx@marginheight#2\relax
  \advance\Mfx@piece@count\@ne
  \MFX@buildpiece@down
  \MFX@buildpiece@up
%    \end{macrocode}
% Once \cs{Mfx@piece@content} has been built, we append it to the
% \cs{Mfx@marginbox}.
%    \begin{macrocode}
  \setbox\Mfx@marginbox\vbox{%
    \unvbox\Mfx@marginbox
    \box\Mfx@piece@content
    \vskip\z@
  }%
}
%    \end{macrocode}
% \end{macros}
%
% \subsection{Third pass: piecewise downward}
% \begin{macros}{\MFX@buildpiece@down,\mfx@pieceout}
% The next pass is another downward pass.  This is very similar
% to the first (global downward) pass, except we're now moving
% the first several notes from \cs{mfx@marginout} into
% \cs{mfx@pieceout} (again, inverting the order) and deferring
% the rest back into \cs{mfx@marginout}.  
%    \begin{macrocode}
\def\MFX@buildpiece@down{%
%<debug>\MFX@debug{buildpiece@down: ENTRY piece=\the\Mfx@piece@count,
%<debug>    marginpos=(\the\Mfx@marginpos@min, \the\Mfx@marginpos@max),
%<debug>    marginspace=\the\Mfx@marginspace,
%<debug>    marginout=\MFX@mac\mfx@marginout}%
  \let\mfx@build@note\MFX@piece@note@down
  \let\mfx@build@skip\MFX@piece@skip@down
  \let\mfx@pieceout\@empty
  \MFX@run@clear\mfx@marginout
%<debug>\MFX@debug{buildpiece@down: RETURN pieceout=\MFX@mac\mfx@pieceout,
%<debug>    marginout=\MFX@mac\mfx@marginout}%
}
%    \end{macrocode}
% \end{macros}
%
% \begin{macros}{\MFX@piece@note@down}
% We again define each of our building macros.  First, the note builder.
% When we encounter a note, we first zero out \cs{Mfx@marginspace}.  Then
% we need to decide whether to put it in the current piece, or whether
% to defer it to the next piece, based on a few signals.  A note will only
% be deferred for one of two reasons.  In either case, we store the decision
% to defer in |@tempswa|, so we'll start by clearing it.
%    \begin{macrocode}
\def\MFX@piece@note@down#1#2#3{%
%<debug>\MFX@debug{piece@note@down: ENTRY: #1at #2, lowest=#3,
%<debug>           marginpos=(\the\Mfx@marginpos@min,\the\Mfx@marginpos@max,
%<debug>           marginspace=\the\Mfx@marginspace}%
  \Mfx@marginspace\z@
  \@tempswafalse
%    \end{macrocode}
% The first reason to defer is if its callout position (|#2|)
% is beneath the end of the current piece (\cs{Mfx@marginheight}) \emph{and}
% if its lowest-possible-piece (|#3|, as computed in the second pass) is
% \emph{after} the current one (\cs{Mfx@piece@count}).
%    \begin{macrocode}
  \ifdim#2>\Mfx@marginheight
    \ifnum#3>\Mfx@piece@count
      \@tempswatrue
    \fi
  \fi
%    \end{macrocode}
% The second possibility is that we've run out of room in this piece.
% In this case, there's no need to check |#3|, since that number was
% computed assuming everything that fit was as low as possible.
%    \begin{macrocode}
  \ifdim\dimexpr\ht#1+\dp#1+\Mfx@marginpos@min>\Mfx@marginheight
    \@tempswatrue
  \fi
%    \end{macrocode}
% If a note is deferred, then we push it and everything after it
% back onto \cs{mfx@marginout}.
%    \begin{macrocode}
  \if@tempswa
%<debug>\MFX@debug{piece@note@down: clearing margin}%
    \MFX@piece@clear
    \mfx@build@note{#1}{#2}{#3}%
  \else
%    \end{macrocode}
% Otherwise, we have decided the note is going into this piece.  In this
% case, we now need to check if any compressible space is needed above it.
% In order to get better alignment for deferred notes, we check for the
% case that the current position is zero and the note's callout position
% is \cs{MFX@minus@inf}.  In this case, we change the callout position
% to instead be the greater of 0 or $\cs{topskip}-\cs{ht}\#1$.
%    \begin{macrocode}
    \dimen@#2\relax
    \ifdim\dimen@=\MFX@minus@inf
      \ifdim\Mfx@marginpos@max=\z@
        \dimen@\topskip
        \advance\dimen@-\ht#1\relax
        \ifdim\dimen@<\z@ \dimen@\z@ \fi
      \fi
    \fi
    \advance\dimen@-\Mfx@marginpos@max
    \ifdim\dimen@>\z@
%<debug>\MFX@debug{piece@note@down: adding compressible \the\dimen@}%
      \MFX@snoc\mfx@pieceout{\noexpand\mfx@build@compressible{\the\dimen@}}%
      \advance\Mfx@marginpos@max\dimen@
    \fi
%    \end{macrocode}
% After (maybe) adding the compressible space, we now add the
% (incompressible) box itself, and accumulate its height in both position
% registers.  We no longer need the piece index or the position, so we
% only store the box itself here.
%    \begin{macrocode}
    \MFX@snoc\mfx@pieceout{\noexpand\mfx@build@note{#1}}%
    \advance\Mfx@marginpos@min\dimexpr\ht#1+\dp#1\relax
    \advance\Mfx@marginpos@max\dimexpr\ht#1+\dp#1\relax
  \fi    
}
%    \end{macrocode}
% \end{macros}
%
% \begin{macros}{\MFX@piece@skip@down}
% Skips are a bit more complicated now.  We no longer gobble the
% initial skips (since the skips at the top and bottom of the page
% have already been eaten).  Instead, we need to look at
% \cs{Mfx@marginspace}: if it's nonzero, we substract it from the
% skip length before adding it incompressibly (if there's any left).
% While we have piece number information here, we just pass it on
% to the upward pass, which can choose to ``late-defer'' initial
% (i.e. the bottom-most) skips in a piece.
%    \begin{macrocode}
\def\MFX@piece@skip@down#1#2{%
  \dimen@#1\relax
  \ifdim\Mfx@marginspace>\z@
    \advance\dimen@-\Mfx@marginspace
    \ifdim\dimen@<\z@ \dimen@\z@ \fi
    \advance\Mfx@marginspace-\dimen@
  \fi
%    \end{macrocode}
% At this point, \cs{dimen@} now stores any further incompressible
% skip we need to add, which is now relative to the top of this piece.
% In that case (note that \cs{Mfx@marginspace} is now necessarily zero),
% we need to check that it actually fits in the current piece, and if
% not, defer it.
%    \begin{macrocode}
  \ifdim\dimen@>\z@
    \ifdim\dimexpr#1+\Mfx@marginpos@min>\Mfx@marginheight
      \MFX@piece@clear
      \mfx@build@skip{\the\dimen@}{#2}%
    \else
%    \end{macrocode}
% If the skip \emph{does} fit, we need to add it to \cs{mfx@pieceout}
% and advance the position registers.
%    \begin{macrocode}
      \MFX@snoc\mfx@pieceout{\noexpand\mfx@build@skip{\the\dimen@}{#2}}%
      \advance\Mfx@marginpos@min\dimen@
      \advance\Mfx@marginpos@max\dimen@
    \fi
  \fi
}
%    \end{macrocode}
% \end{macros}
%
% \begin{macros}{\MFX@piece@clear}
% Finally, we need to handle the case of deferring material.  By
% analogy with the previous two passes, we'll continue to refer to
% this as clearing.  In this case, we need to redefine the note and
% skip macros to save themselves back to \cs{mfx@marginout}.
%    \begin{macrocode}
\def\MFX@piece@clear{%
%<debug>\MFX@debug{piece@clear}%
  \def\mfx@build@note##1##2##3{%
    \MFX@cons\mfx@marginout{\noexpand\mfx@build@note##1{##2}{##3}}}%
  \def\mfx@build@skip##1##2{%
    \MFX@cons\mfx@marginout{\noexpand\mfx@build@skip{##1}{##2}}}%
}
%    \end{macrocode}
% \end{macros}
%
% \subsection{Fourth pass: piecewise upward}
% \begin{macros}{\MFX@buildpiece@up}
% We are now ready for the final pass, where we take the reversed
% list of notes for this piece and stack them up, compressing as much
% compressible space as necessary to make them all fit.  Since this
% is the last pass, we can finally prepend all the boxes into
% \cs{Mfx@piece@content}.  Because it is still possible for pieces
% to have skips at the bottom, we need to worry about which piece
% these skips belong in.  Since we've passed forward the lowest-piece
% information from pass 2, we can choose at the last possible moment
% to defer the bottom-most skips of a piece to the next piece (by
% prepending it to \cs{mfx@marginout}.  Note that compressible spaces
% will \emph{never} be at the beginning of \cs{mfx@pieceout}.  We
% store the excess space in \cs{Mfx@marginheight}.
%    \begin{macrocode}
\def\MFX@buildpiece@up{%
  \Mfx@marginheight\dimexpr\Mfx@marginpos@max-\Mfx@marginheight\relax
  \ifdim\Mfx@marginheight<\z@\Mfx@marginheight\z@\fi
%<debug>\MFX@debug{buildpiece@up: excess=\the\Mfx@marginheight}%
  \let\mfx@build@note\MFX@piece@note@up
  \let\mfx@build@compressible\MFX@piece@compressible@up
  \let\mfx@build@skip\MFX@piece@skip@maybedefer
  \MFX@run@clear\mfx@pieceout\relax
}
%    \end{macrocode}
% \end{macros}
%
% \begin{macros}{\MFX@piece@skip@maybedefer}
% As mentioned earlier, the initial skips in a piece may be deferred
% to later pieces, provided they can possibly go there.  Here we
% check the |#2| argument against \cs{Mfx@piece@count} and defer if
% it is larger.  Otherwise, we switch back to the standard behavior
% defined in \cs{MFX@piece@skip@up}.  Deferred skips do not need to
% be subtracted from the excess because they were never compressible
% in the first place.
%    \begin{macrocode}
\def\MFX@piece@skip@maybedefer#1#2{%
  \ifnum#2>\Mfx@piece@count
%<debug>\MFX@debug{piece@skip deferring: #1, #2 (\the\Mfx@piece@count)}%
    \MFX@snoc\mfx@marginout{\noexpand\mfx@build@skip{#1}{#2}}%
  \else
    \let\mfx@build@skip\MFX@piece@skip@up
    \mfx@build@skip{#1}{#2}%
  \fi
}
%    \end{macrocode}
% \end{macros}
%
% \begin{macros}{\MFX@piece@note@up}
% Now that we've taken care of late deferrals, we can define the
% standard behavior without worrying as much about that.
% These macros finally set their contents into \cs{Mfx@piece@content},
% as well as reset \cs{mfx@build@skip} back to its standard value.
%    \begin{macrocode}
\def\MFX@piece@note@up#1{%
%<debug>\MFX@debug{piece@note@up: #1}%
  \setbox\Mfx@piece@content\vbox{%
    \box#1%
    \unvbox\Mfx@piece@content}%
  \let\mfx@build@skip\MFX@piece@skip@up
}
%    \end{macrocode}
% \end{macros}
% \begin{macros}{\MFX@piece@skip@up}
% Skips are also straightforward.
%    \begin{macrocode}
\def\MFX@piece@skip@up#1#2{%
%<debug>\MFX@debug{skip@up: #1=\the\dimexpr#1\relax (#2)}%
  \setbox\Mfx@piece@content\vbox{%
    \vskip#1\relax
    \unvbox\Mfx@piece@content}%
}
%    \end{macrocode}
% \end{macros}
%
% \begin{macros}{\MFX@piece@compressible@up}
% Finally we come to the compressible space.  Here, as long as
% there's excess space (\cs{Mfx@marginheight}) we drop the
% compressible space.  Once the excess is exhausted, we insert
% them as normal skips.
%    \begin{macrocode}
\def\MFX@piece@compressible@up#1{%
%<debug>\MFX@debug{compressible@up: #1, excess=\the\Mfx@marginheight}%
  \advance\Mfx@marginheight-#1\relax
  \ifdim\Mfx@marginheight<\z@
    \MFX@piece@skip@up{-\Mfx@marginheight}\relax
    \Mfx@marginheight\z@
  \fi
}
%    \end{macrocode}
% \end{macros}
%
% \section{Cleaning up}
% We need to worry about a few more things.  First, what happens
% if we reach the end of the document and there are still deferred
% margin notes?  We need to be able to dump all the margin notes
% whenever the user wants (i.e. before a new chapter), so we'll
% make a macro \cs{dumpmargins} to do this, and then make sure it
% gets called \cs{AtEndDocument}.  Since we're looping to do this,
% we need to make darned sure that every \cs{newpage} shrinks the
% marginlist.
% \begin{macros}{\dumpmargins}
%    \begin{macrocode}
\def\dumpmargins{%
%<debug>\MFX@debug{dumpmargins}%
  \mfx@flush@deferred
  \loop
  \unless\ifx\mfx@marginlist\@empty
%<debug>\MFX@debug{dumpmargins: marginlist=\MFX@mac\mfx@marginlist}%
    \let\temp@\mfx@marginlist
    \vbox{}\clearpage
    \ifx\temp@\mfx@marginlist
      \PackageError{marginfix}{lost some margin notes%
      \ifx\mfx@marginstart\relax\ (missing \noexpand\unblockmargin)\fi
%<debug>: \MFX@mac\mfx@marginlist
      }\@eha
      \let\mfx@marginlist\@empty % be nicer by just dropping one?
      % TODO: also, set an emergency mode to allow oversized notes
    \fi
  \repeat
}
\AtEndDocument{\dumpmargins}
%    \end{macrocode}
% \end{macros}
%
% \section{User macros}
%
% \begin{macros}{\marginskip}
% Inserting a skip in the margin list is simple.  We need only
% append \cs{mfx@build@skip} to \cs{mfx@marginlist}.
%    \begin{macrocode}
\def\marginskip#1{%
  \MFX@cons\mfx@marginlist{\noexpand\mfx@build@skip{#1}}%
}
%    \end{macrocode}
% \end{macros}
%
% \begin{macros}{\clearmargin}
% Likewise, \cs{clearmargin} is easy too.
%    \begin{macrocode}
\def\clearmargin{%
  \MFX@cons\mfx@marginlist{\noexpand\mfx@build@clear}%
}
%    \end{macrocode}
% \end{macros}
% \begin{macros}{\softclearmargin}
% While we call \cs{softclearmargin} a ``clear margin'', it's
% actually just a big \cs{marginskip}.  This allows us to stack
% multiple copies without backing them all up.
%    \begin{macrocode}
\def\softclearmargin{%
  \marginskip{\the\textheight}%
}
%    \end{macrocode}
% \end{macros}
%
% \begin{macros}{\extendmargin}
% We overload \cs{Mfx@marginheight} to be the amount of extension
% at all times except shipout-time.
%    \begin{macrocode}
\def\extendmargin#1{%
  \advance\Mfx@marginheight#1\relax
}
%    \end{macrocode}
% \end{macros}
%
% \begin{macros}{\mparshift}
% This is as simple as setting the dimen register.  We advance so that
% the shifts are cumulative, but there's not really any point either way.
%    \begin{macrocode}
\def\mparshift#1{%
  \advance\Mfx@mparshift#1\relax
}
%    \end{macrocode}
% \end{macros}
%
% \begin{macros}{\blockmargin,\MFX@blockmargin}
% We need two macros to process the optional bracket argument.
% Essentially, \cs{blockmargin} just checks that
% \cs{mfx@marginstart} is defined and then appends a new item to
% the \cs{mfx@marginpieces} list to indicate a piece that starts
% at \cs{mfx@marginstart} and ends at the current position, potentially
% modified by the argument.  It then resets \cs{mfx@marginstart} to
% \cs{relax} and optionally adds a \cs{softclearmargin} if there was no
% star, to keep margin material separated on either side of the phantom.
%    \begin{macrocode}
\def\blockmargin{%
  \@ifnextchar[%]
    \MFX@blockmargin
    {\MFX@blockmargin[0\p@]}%
}
\def\MFX@blockmargin[#1]{%
  \MFX@inject{%
    \ifx\mfx@marginstart\relax
      \PackageError{marginfix}{two \\blockmargin with no \\unblockmargin}\@eha
    \else
      \MFX@cons\mfx@marginpieces{\noexpand
        \@elt{\mfx@marginstart}{\expandafter\dimexpr\mfx@ypos+#1\relax}}%
      \global\let\mfx@marginstart\relax
      \global\advance\Mfx@piece@count\@ne
    \fi
  }%
}
%    \end{macrocode}
% \end{macros}
%
% \begin{macros}{\unblockmargin,\MFX@unblockmargin}
% This is the other half of \cs{blockmargin}.  It ensures that the
% margin is currently blocked, and if so, sets \cs{mfx@marginstart}
% back to a real dimension (the current page position, plus the optional
% modifier).
%    \begin{macrocode}
\def\unblockmargin{%
  \@ifnextchar[%]
    \MFX@unblockmargin
    {\MFX@unblockmargin[0\p@]}%
}
\def\MFX@unblockmargin[#1]{%
  \MFX@inject{%
    \ifx\mfx@marginstart\relax
      \xdef\mfx@marginstart{\dimexpr\mfx@ypos+#1\relax}%
    \else
      \PackageError{marginfix}{\\unblockmargin with no \\blockmargin}\@eha
    \fi
  }%
}
%    \end{macrocode}
% \end{macros}
%
% \begin{macros}{\marginphantom,\MFX@marginphantom}
% The \cs{marginphantom} command is basically just a concatenation of
% \cs{blockmargin} and \cs{unblockmargin}.  We reimplement it from the
% ground up mainly so that the error messages make more sense.
%    \begin{macrocode}
\def\marginphantom{%
  \@ifnextchar[%]
    \MFX@marginphantom
    {\MFX@marginphantom[0\p@]}%
}
\def\MFX@marginphantom[#1]#2{%
  \ifdim#2<\z@\MFX@marginphantom[#1+#2]{-#2}\else
  \MFX@inject{%
    \ifx\mfx@marginstart\relax
      \PackageError{marginfix}{\\marginphantom while margin blocked}\@eha
    \else
      \MFX@cons\mfx@marginpieces{\noexpand
        \@elt{\mfx@marginstart}{\expandafter\dimexpr\mfx@ypos+#1\relax}}%
      \xdef\mfx@marginstart{\dimexpr\mfx@ypos+#1+#2\relax}%
      \global\advance\Mfx@piece@count\@ne
    \fi
  }%
  \fi
}
%    \end{macrocode}
% \end{macros}
%
% \begin{macros}{\marginnotecontinued,\marginnotecontinues,
%     \mfx@contprefix,\mfx@contsuffix}
% The cues of a split note are kept in \cs{mfx@contprefix} (at the start of
% the rest) and \cs{mfx@contsuffix} (at the end of the first part).  They
% are used when a page is built, so we set them globally, like the mode.
%    \begin{macrocode}
\def\mfx@contprefix{(continued)}
\let\mfx@contsuffix\@empty
\def\marginnotecontinued#1{\gdef\mfx@contprefix{#1}}
\def\marginnotecontinues#1{\gdef\mfx@contsuffix{#1}}
%    \end{macrocode}
% \end{macros}
%
% \begin{macros}{\marginnotebesidefloats,\mfx@besidefloats}
% The mode is kept by name, like that of \cs{marginnotepriority}, and
% \cs{MFX@buildmargin} calls |\MFX@beside@|\meta{mode}.
%    \begin{macrocode}
\def\mfx@besidefloats{reserve}
\def\marginnotebesidefloats#1{%
  \@ifundefined{MFX@beside@#1}%
    {\PackageError{marginfix}{Unknown mode `#1' beside floats}%
       {Use reserve or use.}}%
    {\xdef\mfx@besidefloats{#1}}%
}
%    \end{macrocode}
% \end{macros}
%
% \begin{macros}{\marginnotepriority,\mfx@priority}
% The mode is kept by name in \cs{mfx@priority}, and \cs{MFX@buildmargin@down}
% calls |\MFX@down@|\meta{mode}, so a mode is valid if that macro exists.
% We expand the name, since we compare it with \cs{ifx}, and the argument
% may be a macro.  The default is \texttt{warn}.
%    \begin{macrocode}
\def\mfx@priority{warn}
\def\marginnotepriority#1{%
  \@ifundefined{MFX@down@#1}%
    {\PackageError{marginfix}{Unknown margin note priority `#1'}%
       {Use one of warn, defer, split or off.}}%
    {\xdef\mfx@priority{#1}}%
}
%    \end{macrocode}
% \end{macros}
%
% \section{Random scribbles}
%
% Later we'll get fancier with putting notes next to top/bottom
% figures but for now, not so much.
%
% In the future we will support the use of \cs{pdfsavepos} and
% \cs{pdflastypos} for more accurately determining where the callouts
% actually were, which will end up going right around here.  But in
% order to work with older versions of \LaTeX, we still need to
% support the old style of using \cs{@pageht} to figure that out, so for
% now that's all we'll do.
%
% \section{Parting words}
% Finish it up.
%    \begin{macrocode}
\makeatother
%</package>
%    \end{macrocode}
% \makeatother
% \Finale
%
%
% \iffalse
%
% The next line of code prevents DocStrip from adding the
% character table to the generated files(s).
\endinput
% \fi
%
%% \CharacterTable
%%  {Upper-case    \A\B\C\D\E\F\G\H\I\J\K\L\M\N\O\P\Q\R\S\T\U\V\W\X\Y\Z
%%   Lower-case    \a\b\c\d\e\f\g\h\i\j\k\l\m\n\o\p\q\r\s\t\u\v\w\x\y\z
%%   Digits        \0\1\2\3\4\5\6\7\8\9
%%   Exclamation   \!     Double quote  \"     Hash (number) \#
%%   Dollar        \$     Percent       \%     Ampersand     \&
%%   Acute accent  \'     Left paren    \(     Right paren   \)
%%   Asterisk      \*     Plus          \+     Comma         \,
%%   Minus         \-     Point         \.     Solidus       \/
%%   Colon         \:     Semicolon     \;     Less than     \<
%%   Equals        \=     Greater than  \>     Question mark \?
%%   Commercial at \@     Left bracket  \[     Backslash     \\
%%   Right bracket \]     Circumflex    \^     Underscore    \_
%%   Grave accent  \`     Left brace    \{     Vertical bar  \|
%%   Right brace   \}     Tilde         \~}
%%
